\documentclass[11pt,letterpaper]{article}
\usepackage[T1]{fontenc}
\usepackage[utf8]{inputenc}
\usepackage{lmodern,microtype}
\usepackage[margin=.94in,headheight=15pt]{geometry}
\usepackage{amsmath,amssymb,amsthm,mathtools}
\usepackage{booktabs,longtable,tabularx,array,enumitem}
\usepackage{xcolor,fancyhdr,listings,xurl,needspace}
\usepackage[colorlinks=true,linkcolor=blue!45!black,urlcolor=blue!45!black,citecolor=blue!45!black]{hyperref}
\usepackage{bookmark}
\definecolor{ink}{RGB}{24,45,66}
\definecolor{pale}{RGB}{245,247,249}
\hypersetup{pdftitle={Uniform predecessor certificates for intuitionistic PRA},pdfauthor={Research prepared for Vladimir Reshetnikov}}
\setlength{\parindent}{0pt}
\setlength{\parskip}{.3em}
\setlength{\emergencystretch}{3em}
\setlist{itemsep=3pt,topsep=5pt,leftmargin=1.6em}
\pagestyle{fancy}\fancyhf{}
\fancyhead[L]{\small Uniform predecessor certificates for iPRA}
\fancyhead[R]{\small Research report, 10 November 2026}
\fancyfoot[C]{\thepage}
\newtheorem{theorem}{Theorem}[section]
\newtheorem{proposition}[theorem]{Proposition}
\newtheorem{lemma}[theorem]{Lemma}
\newtheorem{corollary}[theorem]{Corollary}
\theoremstyle{definition}\newtheorem{definition}[theorem]{Definition}
\theoremstyle{remark}\newtheorem{remark}[theorem]{Remark}
\newcommand{\ipra}{\mathsf{iPRA}}
\newcommand{\FV}{\operatorname{FV}}
\newcommand{\file}[1]{\expandafter\nolinkurl\expandafter{\detokenize{#1}}}
\lstset{basicstyle=\ttfamily\footnotesize,breaklines=true,columns=fullflexible,
 keepspaces=true,showstringspaces=false,frame=single,rulecolor=\color{black!18},
 backgroundcolor=\color{pale},aboveskip=8pt,belowskip=8pt}

% BEGIN evidence/metrics_macros.tex
\newcommand{\PrimaryTokens}{269}
\newcommand{\PrimaryAST}{307}
\newcommand{\PrimaryBytes}{629}
\newcommand{\PrimaryGzip}{373}

% END

\title{\color{ink}\bfseries Uniform Predecessor Certificates\\for Intuitionistic Primitive Recursive Arithmetic\\[7pt]
\large A complete small checker and a boundary for initializer fusion}
\author{Research prepared for Vladimir Reshetnikov}
\date{10 November 2026}
\begin{document}
\maketitle
\begin{abstract}
We give a complete proof presentation of full first-order intuitionistic
primitive recursive arithmetic, together with an executable Python checker,
both interpretation composites, reproducible whole-program measurements,
and a new obstruction to a broad class of proposed arithmetic reductions.
The numerical and logical basis is inherited from the supplied drafts.
The certificate interface combines constructed predecessor computations
with uniform local computation assumptions, so that a submitted counter
does not need a separate successor recognizer.

The selected complete checker has \PrimaryTokens{} normalized tokens,
\PrimaryAST{} AST nodes, and \PrimaryBytes{} normalized bytes. Its exact
source occupies \PrimaryGzip{} bytes after gzip compression at level 9.
The current finite-corpus minima are 277 tokens, 311 AST nodes, and 672
bytes. These comparisons include every theorem-specific executable
component. The main reduction against the unchanged files comes from a
simpler exit-status rejection convention. Matched-driver controls are
reported explicitly; the source-size result is not a new arithmetic
axiomatization or a significant architectural breakthrough.

The new mathematical result extends the existing bounding-coordinate
obstruction to arbitrary fixed finite register systems and arbitrarily
many fixed probes of one supplied update body. Under an explicit data-flow
restriction, every generated term remains bounded by one input coordinate,
even with all unary primitive-recursive maps and all bounded
primitive-recursive observations freely available. Consequently addition,
unrestricted selection, and injective pairing cannot be term-defined in
that class. This is a restricted obstruction, not an impossibility theorem
for relational or changed-domain bi-interpretations.
\end{abstract}
\begin{center}
\fcolorbox{ink}{pale}{\begin{minipage}{.92\linewidth}
\textbf{Research status.} Full iPRA equivalence is retained, and a complete
checker smaller than the unchanged current programs is supplied.
The requested significant new axiomatic breakthrough has \textbf{not}
been established. The report distinguishes a measured implementation
reduction, inherited equivalence proofs, and the new restricted
nonrepresentability theorem.
\end{minipage}}
\end{center}
\clearpage
\setcounter{tocdepth}{1}\tableofcontents
\clearpage

% BEGIN theory/main.tex
\section{What the present iteration establishes}

The source comparison is pinned to ProveIt commit
\begin{center}\small
\file{0300a290bedee5019c974291850b798378438b2e}.
\end{center}
There are 16 current ZIP archives in the supplied directory. Their byte
strings were verified against the Git blob identities obtained from the
GitHub connector, and the directory was rechecked at this exact revision.
The direct and recursive inventories contain respectively 393 and 2,249
Python member occurrences, representing 297 distinct byte strings.
Small build scripts, measurement tools, and test drivers are excluded from
the complete-checker comparison. Complete checkers included as controls
or references are counted. Every source smaller than any claimed minimum
was inspected and classified; the identities and measurements are supplied
in \file{evidence/}.

The attainable coordinate minima are 277 significant normalized Python
tokens, 311 AST nodes, and 672 normalized UTF-8 bytes. They are attained
by different programs. In particular, the 277-token controls already
inside the scalar-successor and constructed-computation archives must not
be overlooked. A minimum among finitely many submitted implementations
is an observed envelope, not a mathematical lower bound on all checkers.

The present artifact makes three separate contributions.
First, it implements a direct predecessor-packet interface without the
\texttt{list(*counter)} successor recognizer used by the flat-computation
draft. Constructing the successor makes that recognizer unnecessary.
Second, it uses ordinary process exit status for rejection; a successful
run also prints \texttt{ACCEPT}. An uncaught checking or parsing exception
rejects with a nonzero exit status. This removes an otherwise redundant
top-level rejection handler. Third, it proves the single-body obstruction
of Section~\ref{sec:single-body-barrier}.

The first contribution is a synthesis of the constructed-computation,
transient-seed, and flat-computation interfaces~\cite{constructed,transient,flat}.
The initialized scalar iterator, Choice-DI logical rules, endpoint
transport, shared continuation, arithmetic bootstrap, and interpretation
proofs are inherited. We make no claim to have invented those components.
The second contribution is a software-interface simplification, and its
effect is removed in the matched-driver comparison. The third extends an
existing invariant rather than replacing the established arithmetic basis.

\section{The complete axiomatization}
\label{sec:complete-theory}

\subsection{Finite numerical descriptions}

There is one numerical sort. Terms are
\[
                 t::=x\mid 0\mid St\mid I(n,a;u.h).
\]
Only the update $h$ binds $u$. Other variables in $h$ are passive
parameters and stay fixed during the iteration. The two computation
schemas, allowed in every retained context, are
\begin{align}
 I(0,a;u.h)&=a,\tag{I0}\label{up:i0}\\
 I(Sn,a;u.h)&=h[I(n,a;u.h)/u].\tag{IS}\label{up:is}
\end{align}
These are ordinary scalar iterations, whose update receives the preceding
state but not the old counter.

The binder notation has an exact first-order elaboration. Starting with
$0,S$, for each fixed earlier-stage term $h(u,\vec p)$ introduce a symbol
$i_h(n,a,\vec p)$ and the displayed two equations. Expressions and proofs
use only finitely many stages. There is no object variable over function
descriptions, and no universally total program evaluator is added.
Coherence between alternate finite descriptions follows from the internal
uniqueness argument below.

\subsection{Logical syntax and inference rules}

Native formulas have the grammar
\[
            A::=(s=t)\mid Q_yA\mid D_{x,y}A,
\]
where the paired binders $x,y$ are distinct. Contexts are finite reusable
lists. The designated contradiction is $0=S0$.
The following is the complete primitive rule inventory:
\[
 \frac{}{\Gamma\vdash A}\quad(A\in\Gamma)\qquad\text{(Hyp)},
\]
\[
 \frac{\Gamma\vdash s=t\qquad\Gamma\vdash F(s)}
      {\Gamma\vdash F(t)}\qquad\text{(Tr)},
\]
\[
 \frac{\Gamma\vdash Q_yA(y)\qquad\Gamma\vdash A(0)}
      {\Gamma\vdash A(t)}\qquad\text{(QE)},
\]
\[
 \frac{\Gamma\vdash D_{x,y}A(x,y)\qquad
       \Gamma,Q_yA(b,y)\vdash F}
      {\Gamma\vdash F}\qquad\text{(DE)},
\]
\[
 \frac{\Gamma,A(t,0)\vdash A(t,Sz)}
      {\Gamma\vdash D_{x,y}A(x,y)}\qquad\text{(DI-case)},
\]
\[
 \frac{\Gamma,E(t,z)\vdash E(t,Sz)}
      {\Gamma\vdash D_{x,y}E(x,y)}\qquad\text{(DI-progress)}.
\]
In Tr, $F$ is any fixed native formula template. All substitutions avoid
capture. In DE, $b$ is fresh for the retained context, conclusion, and
fixed matrix parameters. In DI, $z$ is fresh for the retained context,
witness $t$, and fixed matrix parameters. Each displayed local assumption
is discharged. The progress matrix $E$ must be a \emph{literal equation};
an equation underneath a binder is not an equation matrix for this rule.

I0 and IS are axiom schemas. There are no additional primitive equality,
successor, induction, or ordinary logical axioms. In particular, compound
progress is not an admissible extra rule: it would supply induction for
arbitrary formulas, exceeding the intended induction restriction.

The equivalent certificate presentation uses a computation theorem as a
local assumption:
\[
 \frac{\Gamma,H\vdash A}{\Gamma\vdash A}\quad\text{(Supply)},
 \qquad H\text{ an instance of I0 or IS}.
\]
This is justified by ordinary-assumption cut. Conversely, Supply followed
by Hyp proves a computation axiom. The checker does not import a library
of theorem schemas beyond the two explicitly implemented computations.

\subsection{Structural admissibility and the zero-predecessor case}

Weakening, exchange, contraction, renaming, numerical substitution, and
ordinary-assumption cut are admissible by induction on finite derivations.
Before insertion or substitution, rename local eigenparameters away from
the inserted terms and retained data. Numerical substitution preserves
the root constructor of an induction matrix, so it preserves the
equation-only restriction. At a cut hypothesis leaf insert the supplied
derivation; every other rule commutes with the construction.

Let $B(p):=I(0,p;v.0)$. I0 gives $B(p)=p$. Tr, using this equation as
both premises with template $w=p$, gives $p=p$. The templates $w=s$,
$s=w$, and $f(s)=f(w)$ then give symmetry, transitivity, and numerical
congruence. These are derivations, not extra equality rules.

\begin{lemma}[Boxed zero suffices]
\label{up:boxing}
I0 and successor computations whose predecessor has a nonempty
recursor representation derive every instance of IS.
\end{lemma}
\begin{proof}
For any term $p$, compute at the composite predecessor $B(p)$:
\[
 I(SB(p),a;u.h)=h[I(B(p),a;u.h)/u].
\]
Transport along $B(p)=p$ with the fixed template
$I(Sw,a;u.h)=h[I(w,a;u.h)/u]$. This yields the desired IS instance.
The argument applies to every domain element and uses no standardness
assumption. In the implementation, the only canonical predecessor that
needs this repair is literal zero; all nonempty canonical predecessors
already select the successor branch.
\end{proof}

\section{Why this is full intuitionistic PRA}
\label{sec:ipra-equivalence}

All metatheorems in this section are inherited and are reproduced in full
in the appendices, with attribution~\cite{choice,flat}. They concern
derivability on the whole numerical domain, with actual equality.

\subsection{Deriving induction, cases, and ordinary logic}

Use an unused existential coordinate in DI-progress, then DE and QE.
This derives contextual equation induction:
\[
 \frac{\Gamma\vdash E(0)\qquad\Gamma,E(z)\vdash E(Sz)}
      {\Gamma\vdash E(t)}.
\]
The same construction with DI-case derives zero/successor cases for
arbitrary native formulas from proofs of their zero and successor
instances. Contexts may contain arbitrary formulas; only an induction
matrix is restricted to an equation.

Define $C(v,a,b)=I(v,a;u.b)$. Its endpoint equations are
$C(0,a,b)=a$ and $C(Sv,a,b)=b$. By external recursion on finite formula
syntax define a gate
\[
 G_v(s=t):=(C(v,0,s)=C(v,0,t)),\qquad
 G_v(D_{x,y}A):=D_{x,y}G_vA.
\]
An auxiliary $Q$ can first be represented by $D$ with an unused existential
coordinate. Then $G_0A$ is provable and $G_{Sv}A$ is interderivable with
$A$. Padding $A$ by $D_{u,y}G_yA$ preserves it.

For equation leaves set
$M_v(s=t,r=w):=(C(v,s,r)=C(v,t,w))$.
For paired formulas align their binders and recurse; if only one input
has a binder, pad the equation input first. The maximum remaining binder
depth decreases. Binder congruence, proved with DI-case, establishes
$M_0(A,B)\dashv\vdash A$ and $M_{Sv}(A,B)\dashv\vdash B$.

Put $1=S0$, $2=S1$, $\delta(v)=C(v,1,0)$, and
$p(z)=I(z,0;u.\delta(u))$. Only $p(1)=1$ and $p(2)=0$ are needed.
With fresh unused coordinates, the full logical interface is
\begin{align*}
 A\to_*B&:=D_{u,z}M_z(A,B),\\
 A\wedge_*B&:=D_{u,z}G_zM_{p(z)}(A,B),\\
 \exists_*x\,A(x)&:=D_{x,y}G_yA(x),\\
 A\vee_*B&:=\exists_*x\,M_x(A,B),\\
 \forall_*x\,A(x)&:=A(0)\wedge_*D_{u,z}G_zA(z),\\
 \bot_*&:=(0=1).
\end{align*}
The appendix proves every introduction and elimination rule. For example,
implication introduction uses the retained zero selector as its assumption;
elimination specializes the resulting closure at $1$. Conjunction
projections specialize at $1,2$. Existential elimination uses a fresh
witness, the zero-gate theorem, and specialization at $1$.

Numerical cases yield
$(x=0)\vee_*\exists_*y(x=Sy)$. Its equality-based elimination is essential
when the target occurs in a retained context: one adds an equality and
transports, instead of silently substituting into the retained context.
Under $0=1$, transport through $C$ yields any numerical equality, and
external formula induction with DI-case gives explosion for every native
formula. No excluded middle for arbitrary quantified formulas is used.

\subsection{Arithmetic strength and the ordinary realization}

Let $T_I$ use ordinary intuitionistic first-order equality logic with
genuine falsity, the staged iterator symbols, I0, IS, equation induction,
and the explicit axioms
\[
 x=0\vee\exists y(x=Sy),\qquad (0=S0)\to\bot.
\]
The falsity bridge must be present: positive numerical equations alone
have a singleton ordinary model. Interpreting native falsity as $0=S0$
does not justify omitting that bridge from the ordinary comparison theory.

The inherited arithmetic construction follows this noncircular order:
predecessor from scalar iteration and equation induction; a characteristic
of actual equality on the entire domain; characteristics of fixed
quantifier-free formulas and hence quantifier-free induction; then pairing
and a compiler for arbitrary fixed primitive-recursive descriptions.
The complete finite predecessor terms and their identities are in
Appendix~\ref{sec:hybrid-bootstrap}. Pairing is not assumed in order to
construct predecessor.

For a target recursion with base $g$ and step $h$, using the proved pairing
$\pi$ and projections $L,H$, the compiler takes
\[
 V(n,\vec p)=I\bigl(n,\pi(0,g^*(\vec p));
 u.\pi(SL(u),h^*(L(u),H(u),\vec p))\bigr),
 \qquad f^*(n,\vec p)=H(V(n,\vec p)).
\]
An equation induction proves $L(V(n,\vec p))=n$; the two projection
identities give the required base and step equations. Passive parameters
remain fixed. The compiler is an external recursion on finite descriptions,
not an internal universal evaluator.

\begin{theorem}[Whole-domain bi-interpretability]
\label{up:biinterpretability}
The native calculus is a faithful presentation of an ordinary theory
bi-interpretable and equiconsistent with full first-order iPRA. Both
numerical domain composites are internally isomorphic to the identity by
the relation $x=y$.
\end{theorem}
\begin{proof}
Interpret numerical iteration in iPRA by primitive recursion ignoring the
counter in the update. Interpret formulas by
\[
 (Q_yA)^\circ=A^\circ(0)\to\forall y A^\circ(y),\quad
 (D_{x,y}A)^\circ=\exists x\bigl(A^\circ(x,0)\to\forall yA^\circ(x,y)\bigr),
\]
and preserve equality. Case DI is valid by coverage and progress DI by
equation induction. All other native rules preserve these readings.
In the reverse direction, the displayed logical macros, equality
reflection, quantifier-free induction, and scalar compiler interpret every
ordinary symbol and axiom of iPRA.

Two iterations with the same seed and provably equal fixed updates are
equal at every counter by equation induction. The same uniqueness argument
applies to general primitive recursion. External induction on a finite
description therefore identifies its double translation with the original
operation. This also proves coherence of the staged binder notation under
substitution. Domains are the whole original domains, and equality is
actual equality, so the graph $x=y$ is internally total, bijective, and
preserves all interpreted operations.

For formulas, the nontrivial native return translation is
\[
 D_{x,y}A\dashv\vdash
 \exists_*x\bigl(A(x,0)\to_*\forall_*y A(x,y)\bigr).
\]
Forward, eliminate D at a fresh witness and use its closure under the
temporarily assumed base. Reverse, eliminate the translated existential
and use its implication and universal instance at $Sz$ to discharge
DI-case. The Q round trip is its unused-witness instance. External induction
on formula syntax gives both full formula composites. The explicit falsity
bridge preserves contradictions; effective finite proof translations give
equiconsistency, without claiming an internal proof of consistency.
\end{proof}

\section{The predecessor certificate interface}
\label{sec:certificate-interface}

The numerical decoder and logical checking operations are inherited from
the flat-computation checker. The changed arithmetic instruction is
\[
                    [\mathtt{null},p,a,h].
\]
It constructs the following local equation from its three data fields:
\[
 \begin{cases}
 I(0,a;u.h)=a,&p\text{ has a false raw representation},\\
 I(Sp,a;u.h)=h[I(p,a;u.h)/u],&p\text{ has a true raw representation}.
 \end{cases}
\]
The condition is on the finite representation, not on whether the value
of the predecessor is zero in a model. A nonempty representation may
denote zero. The source is always explicitly zero or a constructed
successor, so a separate check that a supplied counter has successor
shape is unnecessary.

Canonical literal zero uses the base branch. To obtain the successor step
at zero, use the nonempty predecessor $B(0)$ and
Lemma~\ref{up:boxing}. Every other canonical predecessor, including an
allocated rigid parameter whose value may be zero, selects a valid
successor computation directly. This proves canonical completeness of
the arithmetic interface independently of Python truth values.

\subsection{The shared raw representation}

The public input is one JSON array. Its last element is a requested closed
formula. Earlier four-field instructions are consumed right to left;
context membership closes a goal without reading an instruction. A major
proof is checked before its continuation. An unread prefix is inert.

The same finite tree representation is read in two sorts. Numerically,
an empty list is zero, a unary list is successor, a binary list is zero
with checked annotations, and a ternary list is an iterator. As a formula,
these are respectively truth $0=0$, Q, equality, and D (whose first two
fields are checked annotations). A ternary numerical node gives its
counter and seed two dummy zero slots and its update two aliases of its
one state slot. This does not add a counter argument to the update.

Opening carries a depth $d$, even at formula positions and odd at
numerical positions. Unary, binary, and ternary nodes add respectively
$2,1,4$. At numerical depth $2q+1$, coordinates $2j-1,2j$ name the same
$j$th protected slot; coordinate zero denotes zero. Negative coordinates
name allocated rigid numerical parameters. A global floor starts at zero
and decreases once per continuation. A scalar coordinate is retained only
when it lies between this floor and the current depth. Otherwise a
numerical position uses the supplied image, itself checked at numerical
root, while formula padding completes to zero. Missing used numerical
images reject. The full input/alias contract is in \file{code/FORMAT.txt}.

The raw codec is compact but less legible than the mathematical syntax.
The separately executable readable checker documents the same operations.
It is included in the measurements and is not passed off as the compact
program. Neither implementation imports a theorem library, compiler, or
arithmetic evaluator.

\subsection{Soundness of uniform computation supply}

The only new proof obligation relative to the inherited logical kernel is
the arithmetic source constructor. It is useful to make the treatment of
unused state slots explicit.

\begin{lemma}[Dummy zero slots]
For any fixed successfully traversed raw numerical expression, inserting
dummy zero slots outside its local binders does not change its
zero-completed numerical meaning. Admitted rigid negative parameters may
have arbitrary values.
\end{lemma}
\begin{proof}
Induct on the finite raw tree. An exposed coordinate is completed to zero
in both readings; an internal protected coordinate keeps its local slot
meaning; a rigid parameter keeps its value. Unary numerical nodes add the
same dummy zero, and binary numerical nodes are zero in both readings.
For a ternary node the induction hypothesis identifies the counter and
initializer and identifies the updates at every state. Iterator uniqueness
then identifies the iterations. This is an argument for each fixed finite
expression, not an assumed uniformly total evaluation function.
\end{proof}

\begin{theorem}[Uniform predecessor supply]
\label{up:uniform-supply}
Every successfully supplied computation equation is a theorem of the
native calculus, uniformly in every admitted new parameter. No
known-zero-witness discharge is required.
\end{theorem}
\begin{proof}
For a null major, the common continuation opens the generated equation
with numerical image zero. All newly admitted rigid parameters remain
rigid under this opening.

If the raw predecessor is false, successful traversal makes its numerical
meaning zero. Both occurrences of the initializer agree by the dummy-slot
lemma, and the generated equation is I0.

If the predecessor is true, the constructed source counter is its unary
wrapper, hence a genuine successor regardless of the predecessor's
numerical value. The proposed recursive image is the iterator at that
predecessor, the same initializer, and the same update.
If opening the update uses that image, the image is first checked at the
old parameter floor. The inherited opening/substitution lemma identifies
the resulting right side with the update at that preceding iterator.
Inserting the extra dummy slot for the source successor does not alter
the predecessor's zero-completed meaning. The equation is IS.

If opening never uses the recursive image, the update is independent of
that exposed state coordinate. Any meaningful use of an as-yet
unallocated parameter in the update would enter the same substitution
branch, so such a use is also absent. The predecessor and initializer
may acquire a fresh parameter, but IS with a state-independent update is
valid for every predecessor and initializer. Thus the equation is again
an ordinary computation theorem uniformly in that parameter.

For canonical sources the argument is immediate substitution in I0/IS;
the case distinction above additionally covers successful noncanonical
zero-completed representations. Unlike the older packet, no arbitrary
non-successor term can be mistaken for a successor counter: the wrapper
is constructed by the checker itself.
\end{proof}

\begin{theorem}[Checker correspondence]
Every successful execution proves its decoded closed native goal. Every
closed native theorem has a finite canonical certificate accepted when
the checking computation has sufficient resources.
\end{theorem}
\begin{proof}
Induct on a finite successful execution. Membership is Hyp. A checked
equality or Q major with matching endpoints implements Tr or QE. If an
optional endpoint request does not match, it leaves the old goal to be
proved and adds only truth. A checked D major opens its Q clause at a
fresh rigid parameter. The two introduction choices are exactly DI-case
and literal-equation DI-progress. The decreasing floor validates the
retained goal and witness before allocating a new parameter. Uniform
computation supply is Theorem~\ref{up:uniform-supply}. A malformed scalar
formula has no finite closing rule that creates a new substantive formula:
context entries are the traversed formula lists generated by these cases.
These facts give a finite native derivation by ordinary cut.

For completeness, rename every local binder and eigenparameter in a
native derivation. Use the inherited named encoding of numerical and
formula syntax, emit its logical rule instructions, and replace each
computation axiom by Supply followed by Hyp. A literal-zero predecessor
in IS is replaced by the boxed proof of Lemma~\ref{up:boxing}. All other
canonical computation sources are direct. Reverse the instruction stream
and append the closed goal. Ordinary structural induction on the derivation
establishes acceptance; the supplied producer exercises this construction
on explicit finite examples and is not a universally verified compiler.
\end{proof}

\subsection{A measurable benefit in proof certificates}

The simpler source constructor has a concrete benefit even though its
matched kernel is one byte larger. In the old flat interface, a scalar
parameter predecessor was not list-shaped, so the delivered encoder used
three command records: $B(p)=p$, computation at $B(p)$, and transport.
The new interface checks computation at that parameter directly in one
record. The arbitrary parameter need not denote a positive number; the
constructed \emph{source counter} is $Sp$ in either case.

The same named derivations, emitted by the frozen old and new producers
and accepted by their respective complete checkers, give:
\begin{center}\small
\begin{tabular}{lrr}
\toprule
Named derivation & Old records & New records\\
\midrule
Nonvacuous zero witness &4&2\\
Atomic progress closure &6&4\\
Arbitrary-term specialization &8&6\\
Contextual induction with seed congruence &48&40\\
\bottomrule
\end{tabular}
\end{center}
These are reproducible counts for the supplied encoders, not lower bounds
on the shortest proofs in either calculus. The 48-to-40 example saves
about $16.7\%$ of the command records. The literal-zero predecessor still
uses three records in both delivered encodings. The complete 14-example
comparison includes serialized byte counts as well as command counts.

Physical Python stack and memory limits can reject very large valid
certificates. This is a resource limitation of the implementation, not an
extra object-theory axiom. Canonical well-founded proofs do not depend on
divergent malformed-input aliases. The selected checker retains explicit
list dispatch rather than using recursion-limit-dependent string decoding.

\section{Source measurements, gzip, and matched controls}
\label{sec:measurements}

All measurements count the whole standalone module: the JSON import,
opening and validation, numerical computations, logical inference,
parameter discipline, input driver, and success output. Python and its
standard JSON library are shared runtime infrastructure. No
theorem-specific executable dependency is omitted. Docstrings and comments
are removed, \texttt{ast.unparse} is applied, and a final newline is added.
Significant tokens, every AST node, and UTF-8 bytes are then counted with
CPython 3.12.14, as in the supplied corpus.


% BEGIN evidence/metrics_table.tex
\begin{center}\small
\begin{tabular}{lrrrr}
\toprule
Complete program & Tokens & AST & Bytes & Gzip-9\\
\midrule
Current token leader & 277 & 317 & 684 & 403\\
Current AST leader & 280 & 311 & 705 & 393\\
Current byte leader & 281 & 322 & 672 & 397\\
Current flat computation & 278 & 314 & 679 & 402\\
Current predecessor joint & 280 & 314 & 684 & 390\\
New predecessor, two verdicts & 278 & 315 & 680 & 399\\
Selected complete checker & 269 & 307 & 629 & 373\\
Readable complete checker & 312 & 359 & 1531 & 1024\\
Ancillary byte profile & 269 & 311 & 616 & 366\\
Ancillary AST profile & 268 & 300 & 649 & 367\\
\bottomrule
\end{tabular}
\end{center}

The same rejection convention applied to the old programs gives:

\begin{center}\small
\begin{tabular}{lrrrr}
\toprule
Complete program & Tokens & AST & Bytes & Gzip-9\\
\midrule
Current token leader (matched driver) & 268 & 309 & 633 & 379\\
Current AST leader (matched driver) & 271 & 303 & 654 & 374\\
Current byte leader (matched driver) & 272 & 314 & 621 & 371\\
Current flat computation (matched driver) & 269 & 306 & 628 & 374\\
Current predecessor joint (matched driver) & 271 & 306 & 637 & 366\\
\bottomrule
\end{tabular}
\end{center}

% END


The Bytes column is normalized-source size; Gzip-9 is the size of the
exact raw source, including any comments in the readable version. The
ancillary profiles combine previously published Python representation
choices with the new packet. They are included for comparison, not
selected as the clearest implementation. The AST profile accepts extra
empty-object aliases and rejects nonempty strings through Python's finite
recursion limit; it could diverge on such strings in an ideal unbounded
interpreter. The primary checker has explicit list dispatch and avoids
that particular limitation.

Gzip is an informational metric, not an optimization objective. The exact
procedure is
\begin{lstlisting}
gzip.compress(source_bytes, compresslevel=9, mtime=0)
\end{lstlisting}
The full gzip member, including header and trailer, is counted. Exact raw
and normalized-source gzip sizes and hashes are in
\file{evidence/checker_metrics.json}. The primary compact source is already
normalized, so those two sizes coincide. The delivered \file{.py.gz}
decompresses exactly to its measured Python source. The runtime and zlib
version are recorded for reproducibility.

The exit-status convention is part of the comparison, not hidden in it.
Removing \texttt{except: print('REJECT'); raise} preserves the accepted
inputs: every exception still exits nonzero, and every success still
prints \texttt{ACCEPT}. The same modification is applied to each frozen
leading comparator in the matched table. Therefore the below-corpus
whole-file result does not imply a smaller mathematical basis. The
predecessor constructor also replaces an unusual eager unpacking check
with direct source construction; its exact source-cost difference is
shown rather than described as a large gain.

\subsection{The entire selected executable}

For inspection, this is the whole measured file. There are no omitted
theorem-specific imports or generated axiom tables.
\begin{lstlisting}[language=Python]
from json import load
k = 0
i = lambda a, t=None, d=0: [i(b, t, d + b' \x02\x01\x04'[len(a)]) for b in a] if isinstance(a, list) else a if k <= a < d else d % 2 and i(t, d=1)

def c(a, g):
    global k
    if i(a) in g:
        return
    m, s, u, t = q()
    if m:
        c(m, g)
        if m in ([s, u], not s and [t]) and i(t, u) == a:
            a = i(t, s)
        h = m[2:]
    elif m is None:
        h = [[s and [s], u, t], i(t, [s, u, t], 1) if s else u]
    else:
        h = a = i(a[2], t, 2)
    k -= 1
    c(i(a, [k]), g + [i(h, k if m or len(a) % 2 == m else 0)])
q = load(open(0)).pop
c(q(), [])
print('ACCEPT')
\end{lstlisting}

\section{Validation and reproducibility}


% BEGIN evidence/validation_summary.tex
All 248 complete-program expectations passed (31 cases, four implementations, ordinary and optimized Python). The suite also passed 24 malformed-input runs and three targeted mutation controls. Both old and new encodings were executed in 28 certificate-comparison runs. The independent arithmetic diagnostic attempted 27,600 packets, generated 15,168 equations, and checked 60,672 parameter valuations; it found zero counterexamples and recorded 0 evaluation-budget skips.

% END


The proof examples exercise computation with nonzero seeds, actual state
use, passive parameters, equality transport under binders, Q and D rules,
both introduction choices, and zero-predecessor repair. Negative cases
exercise false arithmetic conclusions, scope leaks, missing images,
malformed input, and an attempted compound-progress inference.
Boundary cases whose intended arithmetic meaning changed with the new
predecessor packet are reclassified by that meaning; an old malformed
counter is not automatically a malformed predecessor.

The numerical falsification tests evaluate the generated equations under
an independent standard-model decoder. They can find a bad transcription
or counterexample but cannot establish a whole-domain interpretation.
The mathematical conclusions rest on the written proofs, including the
attributed appendices. There is no Lean/Rocq formalization or external
human peer review in this package. Independent assistant reviews are
preserved as reviews, not represented as machine-checked metatheorems.

To reproduce the delivered evidence from a fresh extraction:
\begin{lstlisting}
python3 reproduce.py
python3 code/checker.py < tests/fixtures/I0_arbitrary_seed.json
python3 build.py
\end{lstlisting}
The first command verifies file identities, recomputes the whole-program
and gzip metrics, and runs the delivered finite checks without changing
the payload. The second checks a single certificate. The third rebuilds
the modular and standalone LaTeX report with an installed PDFLaTeX
distribution and no shell escape. These operations need no network.


% BEGIN theory/single_body_barrier.tex
% New mathematical result, prepared for the 2026-11-10 iPRA iteration.
% Inclusion fragment: requires amsmath, amssymb, amsthm and theorem/proposition/corollary environments.
\section{A barrier for finite-state, single-body seed fusion}
\label{sec:single-body-barrier}

The current tied-seed obstruction~\cite{transient,counter} concerns the particular scalar recurrence
$B_h(0)=1+h(0,0)$ and $B_h(m+1)=h(m,B_h(m))$.  The following extension
shows that the obstacle is not removed by finitely many initialization
probes, finitely many internal registers, counter-dependent probe locations,
or arbitrary fixed primitive-recursive processing of the probe results.
The essential restriction is that all passive data and all retained state
reach the seed and update combiners only through evaluations of
\emph{one fixed body}.  This is a statement about numerical terms in the
standard model; it is not an impossibility theorem for relational or
changed-domain interpretations.

\subsection{The exact class of constructors}

For a function $f:\mathbb N^r\to\mathbb N$, call $i$ a bounding coordinate
if some nondecreasing unary primitive-recursive function $H$ satisfies
\[
 f(x_1,\ldots,x_r)\le H(x_i)
 \quad\text{for every }(x_1,\ldots,x_r).
 \tag{BC}\label{sb:bc}
\]
A globally bounded function is admitted separately, including at arity zero.

Start with variables and zero.  One may additionally admit \emph{every}
unary primitive-recursive operation and \emph{every} globally bounded
primitive-recursive operation of arbitrary finite arity.  More generally,
any primitive operation already satisfying \eqref{sb:bc} may be admitted.
Ordinary finite term composition is allowed.

A single-body constructor is specified by fixed finite dimensions and
fixed primitive-recursive maps.  Its only supplied subdescriptions are a
counter term $t$ and one body $h(\vec u,\vec p)$, where $\vec u$ is a
fixed finite tuple of local variables and $\vec p$ lists the passive
variables.  Write $n=t(\vec p)$.  For fixed integers $a,b,d,q\ge1$, take
\[
 \begin{aligned}
 A_i&:\mathbb N\longrightarrow\mathbb N^q &&(1\le i\le a),\\
 U_j&:\mathbb N^{d+2}\longrightarrow\mathbb N^q &&(1\le j\le b),\\
 F_0&:\mathbb N^a\longrightarrow\mathbb N^d,\qquad
 F_1:\mathbb N^b\longrightarrow\mathbb N^d,\qquad
 G:\mathbb N^d\longrightarrow\mathbb N.
 \end{aligned}
\]
The constructor denotes the value $G(\vec r_n)$ determined by
\begin{align}
 \vec r_0&=F_0\bigl(h(A_1(n),\vec p),\ldots,h(A_a(n),\vec p)\bigr),
 \label{sb:seed}\\
 \vec r_{m+1}&=F_1\bigl(
 h(U_1(n,m,\vec r_m),\vec p),\ldots,
 h(U_b(n,m,\vec r_m),\vec p)\bigr).
 \label{sb:step}
\end{align}
The maps may differ between constructor schemes, but each is fixed before
its body is supplied.  They are not additional object-language input
functions.  No passive parameter enters $F_0,F_1,G$ directly: all such
information must pass through an evaluation of the same body $h$.
The target counter $n$ and elapsed counter $m$ may enter the probe locations
$A_i,U_j$, exactly as displayed.  The readout receives the final registers,
not an additional independent data argument.

Every supplied body is earlier finite syntax; nesting constructors is
allowed.  Thus each fixed description has a primitive-recursive standard
interpretation, by structural induction followed by primitive recursion on
$m$.  No uniform evaluator for all primitive-recursive descriptions is
postulated.

The previously studied tied-seed constructor is the special case
$d=1$, $q=2$, $a=b=1$, $A_1(n)=(0,0)$,
$U_1(n,m,r)=(m,r)$, $F_0(z)=z+1$, $F_1(z)=z$, and $G(r)=r$.
The present class also allows many probes at distinct locations, register
vectors, and seed probes whose locations depend arbitrarily on the target
counter.

\begin{theorem}[Single-body information barrier]
\label{thm:single-body-barrier}
Every finite term generated by the preceding primitives, composition, and
single-body constructors is globally bounded or has a bounding coordinate
in the sense of \eqref{sb:bc}. Given bounding-coordinate or global-bound
witnesses for the primitive operations in its finite description, the
proof effectively constructs a suitable coordinate and a nondecreasing
primitive-recursive bound. The existence statement does not require these
witnesses to be supplied as additional language data.
\end{theorem}

\begin{proof}
Induct on finite descriptions, treating a body's local and passive
variables as independent inputs when applying the induction hypothesis.
The primitive cases hold by assumption.  Composition preserves the
property: if the outer function is bounded in its $j$th input, compose
its nondecreasing bound with the bound for the $j$th actual argument.
A globally bounded actual argument gives a constant bound.

It remains to analyze \eqref{sb:seed}--\eqref{sb:step}.  For any fixed
primitive-recursive map with finitely many input and output coordinates,
a nondecreasing unary primitive-recursive envelope bounds all its output
coordinates whenever all its input coordinates are at most the envelope's
argument.  Such an envelope is obtained by taking a maximum over a finite
box.  Fix envelopes $f_0,f_1,g$ for $F_0,F_1,G$, respectively.

If $h$ is globally bounded by $C$, every probe value is at most $C$.
Every register, at every stage, is therefore at most
$\max(f_0(C),f_1(C))$.  Applying $g$ gives a constant bound for the
result, independently of the counter and passive parameters.

Suppose next that a passive variable $p_i$ bounds $h$:
\[
       h(\vec u,\vec p)\le H(p_i)
       \quad\text{uniformly in all }\vec u,\vec p.
\]
All seed and update probes obey the same bound, even though their local
arguments depend on the counters and registers.  Hence the final result
is at most
\[
       g\bigl(\max(f_0(H(p_i)),f_1(H(p_i)))\bigr).
\]
This is a nondecreasing primitive-recursive bound depending on the one
passive coordinate $p_i$ alone.

The remaining case is that one local coordinate $u_\ell$ bounds $h$:
\[
       h(\vec u,\vec p)\le H(u_\ell)
       \quad\text{uniformly in the passive parameters.}
\]
Choose nondecreasing primitive-recursive envelopes $a$ and $u$ such that
all coordinates of every $A_i(n)$ are at most $a(n)$, and all coordinates
of every $U_j(n,m,\vec r)$ are at most $u(s)$ whenever
$n,m,r_1,\ldots,r_d\le s$.  Define by primitive recursion
\[
 \begin{aligned}
 v(0,n)&=f_0(H(a(n))),\\
 v(m+1,n)&=\max\left(v(m,n),
       f_1\bigl(H(u(n+m+v(m,n)))\bigr)\right).
 \end{aligned}
\]
Induction on $m$ proves that every coordinate of $\vec r_m$ is at most
$v(m,n)$.  At the base, all local probe arguments are bounded by $a(n)$.
At the step, the nonnegative sum $n+m+v(m,n)$ bounds the two counters and
every register.  Apply the bounds for $U_j$, then $H$, then $F_1$.
Consequently the result is at most $g(v(n,n))$.  Its monotone envelope
\[
                 K(n)=\max_{z\le n}g(v(z,z))
\]
is unary primitive recursive and independent of every passive parameter.

Apply the induction hypothesis to the actual counter term $t(\vec p)$.
If it is bounded by a constant $c$, the whole constructor is bounded by
$K(c)$.  If $t(\vec p)\le N(p_i)$ for a bounding coordinate $p_i$, then
its result is at most $K(N(p_i))$.  This is the required one-coordinate
bound.  The cases exhaust the possible induction hypotheses for $h$,
completing the proof.  All envelopes, finite maxima, compositions, and
recursions in the argument are effective operations on fixed
primitive-recursive descriptions, relative to the supplied primitive bound
witnesses. A bound for an arbitrary globally bounded PR function cannot in
general be computed from its bare program description. Unary PR primitives
have effective finite-box envelopes, and Boolean predicates have known
bound one, so those extensions need no additional bound oracle.
\end{proof}

\begin{corollary}[What these repairs cannot recover]
\label{cor:single-body-exclusions}
No finite term in this class represents ordinary addition, an unrestricted
numerical two-branch selector, or an injective map
$\mathbb N^2\longrightarrow\mathbb N$.
This remains true even if all unary primitive-recursive functions and all
bounded primitive-recursive predicates are freely available as primitives.
\end{corollary}
\begin{proof}
Fixing either argument of addition at zero leaves its output unbounded
in the other argument, so addition has no bounding coordinate.
For a selector $C(0,a,b)=a$, $C(Sn,a,b)=b$, fixing the control at zero
leaves $a$ unbounded; fixing $a=0$ and the control at one leaves $b$
unbounded; fixing $b=0$ and the control at zero leaves $a$ unbounded.
Thus none of its three coordinates bounds it.  Finally, for a binary
injection, fixing a supposed bounding coordinate at zero would map an
infinite fiber injectively into a finite bounded subset of $\mathbb N$,
a contradiction.  A global bound fails for the same reasons.
\end{proof}

\subsection{The exact boundary and the useful design lesson}

The theorem is not a prohibition on recursion with finitely many
registers.  Its crucial hypothesis is the data interface in
\eqref{sb:seed}--\eqref{sb:step}.  For example, if a passive initializer
$a$ is permitted to enter a register directly, the one-register loop
\[
             r_0=a,\qquad r_{m+1}=Sr_m,
             \qquad\text{output }r_n
\]
already computes $a+n$, which has no bounding coordinate.  Even while keeping initialization tied to a body, direct register
retention escapes the theorem: take $r_0=h(0,\vec p)$ and
$r_{m+1}=Sr_m$.  For the constant body $h(u,a)=a$, the output is again
$a+n$.  This update is deliberately outside \eqref{sb:step}, because
the old register reaches the update combiner without first passing
through the body.  Such changes may be good designs, but they genuinely
change the data interface rather than merely probe the same body more
elaborately.

The theorem therefore rules out a substantial family of natural repairs
to the tied-seed proposal: adding a fixed number of probe locations,
finite internal registers, arbitrary unary decoders, arbitrary bounded
observations, or more elaborate fixed processing of the probe outputs
cannot by itself yield a term compiler for iPRA's primitive-recursive
signature.  It does not rule out a second independent body, an independent
data input, direct register retention, graph-defined functions, a quotient
or changed domain, or a different logical interpretation.  In particular,
no conclusion about all relational bi-interpretations follows from this
standard-model term obstruction.

% END


\section{What remains open}

The retained numerical basis already has only zero, successor, and one
scalar-iteration constructor. The logical presentation has equality and
one paired binder, with Q a definitional convenience. Its apparent brevity
rests on substantial derivations of predecessor, equality reflection,
ordinary logic, and the primitive-recursion compiler. Counting just a
list of primitive rule names would hide that explanatory burden.

The newer successor-seeded scalar basis removes the separate successor
constructor, but its predecessor and logical bootstrap have already been
completed in the current drafts~\cite{scalar,completion,transient}.
Repeating those proofs would not be a new iteration. The tied-seed
variants have a genuine data-flow barrier; the new theorem extends its
scope but expressly leaves other interpretations open.

Deleting atomic DI-case is also not a valid simplification of the
unrestricted existing formula language. Changing the public language to
gated atomic matrices can evade that particular obstruction, but it
introduces new grammar or decoder obligations and has no demonstrated
complete-checker advantage here. Similarly, a graph-defined selector does
not by itself extend an old recursion schema to terms using that new
function. Such a schema extension needs its own conservativity proof.

An absolute minimum, a significantly smaller mathematical architecture,
and a substantial improvement under a fully matched checker convention
remain open. The appropriate conclusion is the concrete implementation
and restricted theorem proved above, not a claim that the user's broader
research objective has been solved.

\appendix
\clearpage
\section*{Attribution and complete inherited proofs}
\addcontentsline{toc}{section}{Attribution and complete inherited proofs}
The following five sections are reproduced unchanged from the
flat-computation archive at the pinned revision. That archive traces them
to the generated-assumptions, named-specialization, and earlier scalar
developments. Their exact source hashes are in
\file{evidence/inherited_provenance.json}. The upstream MIT No Attribution
license is included at \file{reference/PROVEIT_LICENSE.txt}; mathematical
attribution is retained throughout. The sections supply the complete
written bridge used by Theorem~\ref{up:biinterpretability}.


% BEGIN reference/inherited_choice_basis.tex
% New mathematical result: complementary introduction choices.
% Requires amsmath/amssymb/amsthm and theorem, lemma, proposition environments.
% The inherited arithmetic-to-iPRA bridge must be supplied separately.

\section{Complementary introductions with one local assumption}
\label{sec:choice-introduction}

The old paired introduction retains both the zero instance and an
induction instance when its matrix is an equation. Both together are
unnecessary. A certificate may choose either one, provided the induction
choice is restricted to a literal equation. This changes neither the
formula language nor the meaning of a formula. The omitted combination
is recovered by a short nested introduction--elimination construction.

\subsection{Complete native syntax and rule inventory}

The scalar numerical language is
\[
 t::=x\mid0\mid St\mid I(n,a;u.h).
\]
Only $h$ binds the state variable $u$; $n$ and $a$ do not.
The update may contain fixed numerical parameters. The computation
axioms in every retained context are
\[
 I(0,a;u.h)=a,\qquad
 I(Sn,a;u.h)=h[I(n,a;u.h)/u].
 \tag{I0, IS}\label{eq:choice-computations}
\]
The argument also applies unchanged to the full fixed-description PR
language with its projection, composition, and recursion equations.
No numerical variable ranges over function descriptions, and no uniform
evaluator is postulated.

Native formulas and reusable context clauses are
\[
 A::=(s=t)\mid\mathsf D_{x,y}A,\qquad
 \langle A;y.B(y)\rangle,\quad y\notin\operatorname{FV}(A).
\]
Bound variables are renamed before capture-free substitution.
The two variables $x,y$ in a paired formula are distinct.
An ordinary assumption $F$ is the constant clause
$\langle H;v.F\rangle$, where $v$ is absent from $F$ and $H$ is the
fixed computation theorem $I(0,0;u.0)=0$. Thus ordinary hypothesis
use is a clause Use preceded by this computation proof.

There are computation axioms and exactly the following inference
patterns:
\[
\frac{\Gamma\vdash A}{\Gamma\vdash B(t)}
\quad\mathsf{Use}
\quad(\langle A;y.B(y)\rangle\in\Gamma),
\]
\[
\frac{\Gamma\vdash\mathsf D_{x,y}B(x,y)\qquad
      \Gamma,\langle B(a,0);y.B(a,y)\rangle\vdash F}
     {\Gamma\vdash F}
\quad\mathsf{DE},
\qquad
\frac{\Gamma\vdash s=t\qquad\Gamma\vdash F(s)}
     {\Gamma\vdash F(t)}
\quad\mathsf{Tr},
\]
\[
\frac{\Gamma,B(t,\eta_\varepsilon(z))\vdash B(t,Sz)}
     {\Gamma\vdash\mathsf D_{x,y}B(x,y)}
\quad\mathsf{DI}_{\varepsilon},
\qquad
\eta_0(z)=0,\quad\eta_1(z)=z.
\tag{Choice DI}\label{eq:choice-di}
\]
Here $\varepsilon$ is a certificate choice in $\{0,1\}$, and
$\varepsilon=1$ is permitted only when $B$ has a literal equation
root. An equation beneath a binder is not a literal equation.
Call $\varepsilon=0$ the \emph{case choice} and $\varepsilon=1$
the \emph{progress choice}. Both retain exactly one ordinary
formula assumption. Compound matrices have only the case choice.

In DI the eigenparameter $z$ is fresh for $\Gamma$, the witness
$t$, and every fixed parameter of $B$. In DE the eigenparameter
$a$ is fresh for $\Gamma$, the conclusion $F$, and every fixed
parameter of $B$. In transport, $F(v)$ is any fixed native
formula template. Displayed local assumptions and clauses are
discharged. There are no primitive reflexivity, coverage,
separation, ordinary logical, or additional induction rules.

We call this calculus $\mathcal C$. The original comparison
calculus $\mathcal N$ has identical syntax, computation,
Use, DE, and transport, but its introduction is
\[
\frac{\Gamma,B(t,0),[B(t,z)]_{\mathrm{at}}\vdash B(t,Sz)}
     {\Gamma\vdash\mathsf D_{x,y}B(x,y)}.
\tag{Old DI}\label{eq:choice-old-di}
\]
The bracket supplies $B(t,z)$ exactly at a literal equation.

\subsection{Exact equivalence and the unused-witness construction}

\begin{lemma}[Structural admissibility]
\label{lem:choice-structural}
Weakening, exchange, contraction, renaming, ordinary-assumption cut,
and capture-free numerical substitution are admissible in $\mathcal C$.
\end{lemma}
\begin{proof}
Induct on derivations. Rename local eigenparameters before inserting
proofs or terms. Numerical substitution preserves a formula's root
constructor, so an allowed progress choice remains allowed.
The selected choice does not change under substitution or cut.
For ordinary-assumption cut, replace each use of its constant clause
by the supplied proof. Computation schemas are closed under
capture-free numerical substitution. These facts verify every case.
\end{proof}

\begin{theorem}[Identity-language equivalence]
\label{thm:choice-equivalence}
For every retained clause context $\Gamma$ and native formula $F$,
\[
 \Gamma\vdash_{\mathcal C}F
 \quad\Longleftrightarrow\quad
 \Gamma\vdash_{\mathcal N}F.
\]
The transformations are effective on finite derivations and leave
all numerical and formula syntax unchanged.
\end{theorem}
\begin{proof}
Each new introduction is an instance of the old one after weakening
by the omitted local assumption. For the case choice at an equation,
add the induction instance. For the progress choice, add the zero
instance. Compound case introduction is already the old rule.
Thus every new derivation translates to an old one.

For the reverse direction, transform an old derivation recursively.
Only its atomic DI case needs a construction. Suppose its matrix is
the equation $E(x,y)$ and its transformed premise is
\[
 \pi:\quad\Gamma,E(t,0),E(t,z)\vdash E(t,Sz).
\]
The outer new inference will be case DI with witness $t$. It is
therefore enough to prove $E(t,Sz)$ under
$\Gamma,E(t,0)$.

Rename the eigenvariable in $\pi$ to a fresh $v$. Choose a new
variable $u$ absent from $E(t,y)$ and all retained data. Under
$\Gamma,E(t,0)$, progress DI with witness zero introduces
\[
             \mathsf D_{u,y}E(t,y)
\]
from exactly the renamed premise $\pi[v/z]$.
The bound coordinate $u$ is deliberately unused.

Apply DE to this paired formula. Because $u$ is unused, the opened
clause is exactly
\[
               \langle E(t,0);y.E(t,y)\rangle;
\]
it contains no new existential witness. The retained ordinary
assumption $E(t,0)$ supplies its antecedent. Use at $Sz$ gives
$E(t,Sz)$, and DE closes. Finally the outer case DI introduces
the original $\mathsf D_{x,y}E(x,y)$ with witness $t$.
The original premise proof $\pi$ has been inserted once.

An old compound introduction already is the new case choice.
All other rule cases are unchanged. This completes the induction
on the original derivation.
\end{proof}

The unused inner coordinate is essential. Introducing
$\mathsf D_{x,y}E(x,y)$ internally and then eliminating it would
give an arbitrary fresh witness, not necessarily the original
$t$. The construction instead keeps $t$ as a fixed parameter
of $E(t,y)$ and binds only the unused coordinate $u$.

\begin{proposition}[Linear bound on inference nodes]
\label{prop:choice-node-bound}
Count each computation axiom, Use, DI, DE, and transport as one
inference node, expanding ordinary-assumption use as a computation
of $H$ followed by Use. If an old derivation has $N$ nodes and
$r$ atomic DI nodes, the displayed transformation has exactly
\[
                   N+5r\leq 6N
\]
nodes.
\end{proposition}
\begin{proof}
Each atomic old DI node is replaced by six nodes: outer case DI,
inner progress DI, DE, Use of the opened clause, Use of the
retained base assumption, and the computation of $H$ needed for
that last ordinary-assumption use. Its old premise subtree
occurs once. All other inference nodes are copied once.
\end{proof}

This counts inference-tree nodes. It is not a bound on serialized
JSON bytes, formula sizes, checker runtime, or the number of
consumed packet cells when hypothesis closure is implicit.
In the formula-context Q/D presentation, the corresponding
replacement has outer case DI, inner progress DI, DE, QE,
and two Hyp nodes, giving the same six-node count.

\subsection{Induction and numerical cases separate cleanly}

\begin{lemma}[Contextual equation induction]
\label{lem:choice-ei}
The progress choice, DE, and Use derive
\[
\frac{\Gamma\vdash E(0)\qquad
      \Gamma,E(z)\vdash E(Sz)}
     {\Gamma\vdash E(t)}
\tag{EI}\label{eq:choice-ei}
\]
for every equation $E$, with the usual freshness condition.
\end{lemma}
\begin{proof}
Choose an unused witness coordinate $u$. Progress DI introduces
$\mathsf D_{u,y}E(y)$ from the given step. DE opens the clause
$\langle E(0);y.E(y)\rangle$. Supply its antecedent using the given
base proof and use the clause at $t$. Close DE.
\end{proof}

\begin{lemma}[Cases, equality, and binder congruence]
\label{lem:choice-preinduction}
The case choice derives, for arbitrary native formulas,
\[
\frac{\Gamma\vdash F(0)\qquad\Gamma\vdash F(Sz)}
     {\Gamma\vdash F(t)}
\tag{NC}\label{eq:choice-nc}
\]
with the usual freshness condition. Equality reflexivity,
symmetry, transitivity, numerical congruence, and uniform
matrix interderivability under $\mathsf D$ are derivable
without using the progress choice.
\end{lemma}
\begin{proof}
For cases, use case DI with an unused witness coordinate and
matrix $F(y)$, ignoring its retained base in the successor
premise. DE opens $\langle F(0);y.F(y)\rangle$; the given zero
proof and Use at $t$ finish.

For reflexivity put $r=I(0,t;u.0)$. Use the computation equation
$r=t$ as both transport premises with motive $v=t$ to obtain
$t=t$. The usual explicit equality motives give symmetry,
transitivity, and numerical congruence.

For binder congruence, open the old paired formula at a fresh
witness $a$ and introduce the new one by case DI at $a$.
Convert its retained base to the old base, use the opened
old clause at the successor, and convert back. Ordinary cut
closes the construction, and the reverse conversion gives
the converse. No induction is used.
\end{proof}

The primitive progress choice is still restricted to equations.
Allowing it at arbitrary compound matrices would, by the
argument for \eqref{eq:choice-ei}, supply induction for
arbitrary formulas and change the intended theory.

\subsection{The entire ordinary logical interface}
\label{subsec:choice-logic}

Here is a predecessor-free reconstruction of the ordinary logical
operations. It makes explicit that the case choice supplies the
logical interface independently of progress induction.

Let $C(v,a,b):=I(v,a;u.b)$ with $u$ unused. Computation gives
$C(0,a,b)=a$ and $C(Sv,a,b)=b$. Define the gate on fixed finite
formula syntax by
\[
 G_v(s=t):=(C(v,0,s)=C(v,0,t)),\qquad
 G_v(\mathsf D_{x,y}B):=\mathsf D_{x,y}G_vB,
\]
renaming bound variables away from the control. Equality,
binder congruence, and external structural induction give
\[
 \vdash G_0A,\qquad G_{Sv}A\dashv\vdash A,\qquad
 A\dashv\vdash\mathsf D_{u,y}G_yA
 \quad(u,y\notin\operatorname{FV}(A)).
\tag{Gate and padding}\label{eq:choice-gate}
\]
For padding, case DI proves the forward direction. DE, the
zero-gate theorem, and Use at $1$ prove the reverse direction.

Define a selector between two arbitrary formulas by
\[
 M_v(s=t,u=w):=(C(v,s,u)=C(v,t,w))
\]
and, with aligned fresh paired variables, by
\[
\begin{aligned}
 M_v(\mathsf D_{x,y}A,\mathsf D_{x,y}B)
   &:=\mathsf D_{x,y}M_v(A,B),\\
 M_v(E,\mathsf D_{x,y}B)
   &:=\mathsf D_{x,y}M_v(G_yE,B),\\
 M_v(\mathsf D_{x,y}A,E)
   &:=\mathsf D_{x,y}M_v(A,G_yE),
\end{aligned}
\]
where $E$ is an equation. Every call decreases maximum input
binder depth. Padding and congruence prove
$M_0(A,B)\dashv\vdash A$ and
$M_{Sv}(A,B)\dashv\vdash B$.

Put $\delta(v)=C(v,1,0)$ and $p(z)=I(z,0;u.\delta(u))$.
Only the fixed computations $p(1)=1$ and $p(2)=0$ are needed.
The full logical macro inventory is
\[
\begin{aligned}
 A\Rightarrow_*B&:=\mathsf D_{u,z}M_z(A,B),\\
 A\wedge_*B&:=\mathsf D_{u,z}G_zM_{p(z)}(A,B),\\
 \exists_*x\,A(x)&:=\mathsf D_{x,y}G_yA(x),\\
 A\vee_*B&:=\exists_*x\,M_x(A,B),\\
 \forall_*x\,A(x)&:=
        A(0)\wedge_*\mathsf D_{u,z}G_zA(z),\\
 \bot_*&:=(0=1).
\end{aligned}
\tag{Logical macros}\label{eq:choice-macros}
\]
Newly bound variables are fresh; $u$ is unused.

Implication introduction obtains its assumption from the zero
selector and supplies the consequent at successors, then uses
case DI. Elimination opens its clause and uses it at $1$.
Existential introduction uses its witness and the successor
gate; elimination opens at a fresh witness, supplies the
zero-gate theorem, and uses the clause at $1$.

Disjunction introduction uses witnesses $0$ and $1$.
For elimination, discharge the two branch proofs as
$A\Rightarrow_*F$ and $B\Rightarrow_*F$. NC applied to
$M_z(A,B)\Rightarrow_*F$ gives this implication at every
control: the two endpoint cases follow from the branch
implications. Existential elimination provides $M_a(A,B)$
and the implication gives $F$.

Given $A$ and $B$, selector endpoints and NC give $M_v(A,B)$
at an arbitrary control. Substitute $v=p(Sz)$ and use the
successor gate for conjunction introduction. For projections,
open its zero-gated clause and specialize at $1$ and $2$;
the two fixed computations give $B$ and $A$, respectively.

These rules prove existential numerical coverage:
\[
     (x=0)\vee_*\exists_*v\,(x=Sv).
\tag{Coverage}\label{eq:choice-coverage}
\]
Use NC. Zero takes the left injection of reflexivity, and a
successor takes the right injection with the evident witness.
Coverage gives equality cases even when the retained context
contains $x$; its elimination adds $x=0$ or $x=Sv$ rather
than silently changing that retained context.

Universal introduction follows from a fresh-parameter proof
by substituting zero and successors and forming the two
conjuncts. For elimination, take its $A(0)$ conjunct and
open the paired conjunct. Its base is a zero-gate theorem.
Coverage of the target chooses the zero instance from
$A(0)$ or a successor instance from the opened clause
and gate removal.

Finally, under $0=1$, computation and transport give
$s=C(0,s,t)=C(1,s,t)=t$ for arbitrary numerical terms.
Every equation follows, and external induction on formula
syntax with case DI gives every native formula. This
supplies falsity elimination. No excluded middle for
quantified formulas has entered the construction.

\subsection{Whole-domain realizations and both formula composites}

For a standard first-order reading, let $T_{\rm eq}$ have
intuitionistic equality logic, the same numerical signature
and computation equations, numerical coverage, separation
$(0=1)\to\bot$, and contextual equation induction.
Its interpretation of the native syntax is
\[
 (s=t)^\circ=(s=t),\qquad
 (\mathsf D_{x,y}B)^\circ
   =\exists x\bigl(B^\circ(x,0)\to\forall y B^\circ(x,y)\bigr),
\]
\[
 \langle A;y.B\rangle^\circ
       =A^\circ\to\forall y B^\circ.
\]
The domain is the entire numerical domain and interpreted
equality is actual equality.

Case DI is sound in this realization by numerical coverage.
Progress DI is sound by contextual equation induction under
the conditional base. Use, DE, transport, and computation
have their displayed ordinary meanings. Conversely, the
native macros derive every ordinary logical rule,
\eqref{eq:choice-ei} derives equation induction, and
\eqref{eq:choice-coverage} derives coverage. Separation
translates to the identity implication
$(0=1)\Rightarrow_*(0=1)$.

For the essential paired formula composite,
\[
 \mathsf D_{x,y}B
 \dashv\vdash
 \exists_*x\bigl(B(x,0)\Rightarrow_*\forall_*y\,B(x,y)\bigr).
\tag{Paired composite}\label{eq:choice-paired-composite}
\]
Forwards, DE opens its clause at a fresh $a$; under $B(a,0)$
that clause supplies $B(a,y)$ at a fresh arbitrary $y$.
Introduce the universal, implication, and existential
macros and close DE. Backwards, eliminate the existential
at a fresh $a$. Introduce the paired formula with witness
$a$ by case DI: the retained base, supplied implication,
and universal elimination at $Sz$ give its successor premise.

The ordinary interpretations of the six logical macros
have their ordinary round trips. For conjunction, its true
base and the instances $1,2$ recover both operands;
the converse uses selector cases. The universal has all
successor instances in its paired conjunct and zero in
its first conjunct; coverage gives exactly the ordinary
universal. Implication, disjunction, and existential
quantification use their selector or gate endpoints.
Falsity uses the explicit ordinary separation bridge.
External induction on finite formulas, together with
\eqref{eq:choice-paired-composite}, gives both formula
composites. Clauses correspond in both directions to
$A\Rightarrow_*\forall_*yB$ by the derived logical rules.

The numerical composites are literal identities with
internally provable graph $x=y$. No standard-numeral
restriction, retract, quotient, or equality reflection
assumption is being inserted at this stage. The separately
supplied arithmetic bridge identifies $T_{\rm eq}$ with
full iPRA; the present change preserves that bridge
unchanged because Theorem~\ref{thm:choice-equivalence}
preserves every native sequent.

\subsection{What the theorem does and does not claim}

The new inference interface retains one ordinary assumption
at every paired introduction. It preserves the entire
consequence relation, not just equations or closed
theorems. Consequently the established full-iPRA
interpretations and equiconsistency transfer unchanged.

This result is different from retaining only progress
introduction and strengthening closure elimination.
Here the original elimination rule is retained, and
the case and progress introductions coexist. It is also
different from allowing arbitrary compound progress,
which would supply stronger induction.

A complete-checker comparison must still count the
certificate choice, its equation-root restriction,
substitution, context handling, computation, and the
public input/output interface. The mathematical theorem
alone is not a byte-count or minimality theorem.
The inference-node bound is distinct from source-size
measurements and from serialized certificate size.
The proof is a written universal proof; finite regression
tests do not constitute a Lean/Rocq formalization of it.

% END


% BEGIN reference/inherited_ordinary_bridge.tex
\section{Arithmetic strength and both ordinary interpretation composites}
\label{sec:ordinary-bridge}

We spell out the arithmetic bridge to avoid identifying standard numerical
agreement with bi-interpretability. First elaborate every fixed iterator
into an ordinary function symbol. At stage zero the symbols are $0,S$.
Given a fixed earlier-stage term $h(u,\vec p)$, introduce
$i_h(n,a,\vec p)$ with the two equations
\[
i_h(0,a,\vec p)=a,\qquad
i_h(Sn,a,\vec p)=h(i_h(n,a,\vec p),\vec p).
\]
Only finite descriptions occur; the domain contains no code for a
universal evaluator. Formal parameters are distinct when a description
is declared, while actual arguments can coincide.

Let $T_I$ be ordinary intuitionistic first-order logic with actual
equality and genuine falsity, these staged equations, contextual
induction for literal equations, and
\[
\forall x\,(x=0\vee\exists y\,x=Sy),
\qquad (0=S0)\to\bot.
\]
The second sentence is the falsity bridge. It must be stated: positive
numerical equations by themselves admit a singleton structure. Native
explosion from $0=S0$ and a translation of ordinary falsity to that
equation do not authorize silently dropping the bridge from $T_I$.

The native-to-ordinary translation is
\[
(s=t)^\circ=(s=t),\qquad
(\mathsf D_{x,y}B)^\circ=
\exists x\bigl(B^\circ(x,0)\to\forall y B^\circ(x,y)\bigr).
\]
A clause $\langle A;y.B\rangle$ becomes $A^\circ\to\forall yB^\circ$.
For the auxiliary closure used by the compact checkers,
$(\mathsf Q_yB)^\circ=B^\circ(0)\to\forall yB^\circ$.
The interpretation domain is all elements, defined by $x=x$;
equality is actual equality.

Case DI is sound because numerical coverage combines its retained base
with its successor proof. Progress DI is sound by literal-equation
induction under a temporarily assumed base. Use, DE and transport are
ordinary logical inferences under the translations. Conversely the
derived logical operations, equation induction, cases and explosion in
Section~\ref{sec:choice-introduction} recover the rules and axioms of
$T_I$. An auxiliary closure is interderivable with
$\mathsf D_{v,y}B(y)$ for unused $v$: in one direction use case DI with
the closure and its base; in the other eliminate the paired formula
and return its supplied closure. Thus formulas and context clauses have
the stated meanings in both presentations.

\subsection{The arithmetic dependency order}

The scalar update has no recursion-counter argument. We do not conceal
one in an encoding. The inherited internal proofs, reproduced in the
appendices, establish the required arithmetic in this order:
\begin{enumerate}
\item Using equation induction and numerical cases, construct a fixed
scalar predecessor term and prove $P0=0$ and $P(Sn)=n$.
Existential coverage then gives $x=0\vee x=S(Px)$: in the successor
case $x=Sy$, apply $P$ to obtain $Px=y$ and substitute.
\item Construct truncated subtraction and a numerical characteristic
of actual equality. Prove its reflection on the entire original domain.
\item Compile each fixed quantifier-free formula into its characteristic
equation and recover every contextual quantifier-free induction instance.
\item With that induction available, construct pairing and projections
by a bounded square-root search and compile every fixed general
primitive-recursive description into a scalar iterator term.
\end{enumerate}
No pairing or faithful equality detector is assumed to obtain
predecessor. No appeal to standard-model correctness replaces an
internal equality. Every induction fixes its passive parameters;
diagonal substitution occurs only into a previously proved identity.

For a target recursion
\[
f(0,\vec p)=g(\vec p),\qquad
f(Sn,\vec p)=h(n,f(n,\vec p),\vec p),
\]
let $\pi,L,H$ be the proved pairing and projections. The compiler uses
\[
V(n,\vec p)=I\bigl(n,\pi(0,g^*(\vec p));
u.\pi(SL(u),h^*(L(u),H(u),\vec p))\bigr),
\qquad f^*(n,\vec p)=H(V(n,\vec p)).
\]
Equation induction proves $L(V(n,\vec p))=n$; the two projection
identities then prove the target defining equations. The construction
is external recursion on each finite PR description, not an internal
quantification over arbitrary functions.

\begin{theorem}[Whole-domain bi-interpretability and equiconsistency]
The ordinary realization $T_I$ and full first-order iPRA are
bi-interpretable. In both theories $x=y$ defines the isomorphism from
the original domain to its double interpretation. The native Choice-DI
calculus is therefore a faithful proof presentation of a theory
bi-interpretable and equiconsistent with iPRA.
\end{theorem}
\begin{proof}
Interpret an iterator in iPRA by primitive recursion that ignores its
counter argument in the update. iPRA proves its equations, literal
equation induction, coverage and separation. In the reverse direction
use the scalar compiler. The predecessor and equality-reflection proofs
supply the successor axioms and every quantifier-free induction instance.
All ordinary logical rules are preserved.

For two iterators with the same seed and provably equal updates,
equation induction on the equality of their values proves uniqueness at
every counter, with all parameters fixed. The corresponding argument
applies to ordinary primitive recursion. Induct externally on a finite
description: its double translation satisfies its original defining
equations, so internal uniqueness identifies the two functions. This
also proves coherence when substitution changes the staged syntactic
description of an update.

The domain and equality are unchanged; hence the identity graph is
internally total, bijective and preserves every interpreted operation.
The displayed ordinary meanings of the logical macros and the paired
binder give both formula composites by external induction on finite
formula syntax. Clauses are similarly equivalent to their conditional
universal meanings. The explicit falsity bridge preserves contradiction
in both directions. Effective translations of finite proofs give
equiconsistency, without a claim that either theory proves its own
consistency.
\end{proof}

% END


% BEGIN reference/scalar_predecessor.tex
% Adapted with attribution from the pinned unparenthesized one-hole archive.
\section{Internal scalar predecessor bootstrap}
\label{sec:hybrid-bootstrap}

This scalar bootstrap is inherited from the supplied scalar development~\cite{erasure} and is reproduced to keep the scalar checker self-contained. It is not a new contribution of this paper. The proof uses a factor map fixing odd states and a saturated-iteration construction of predecessor.

The available arithmetic terms are $0$, $S$, and
\[
 I(n,a;u.h),\qquad
 I(0,a;u.h)=a,\qquad
 I(Sn,a;u.h)=h[I(n,a;u.h)/u].
\]
The update binds only $u$.  All other variables in it are fixed numerical
parameters.  We use equality reasoning, the already reconstructed
zero/successor cases, and induction on a single literal equation.
Every induction below keeps its other parameters fixed.  Predecessor,
successor injectivity, pairing, and induction on arbitrary
quantifier-free formulas are not assumptions of the argument.

\subsection{Elementary terms and two finite cycles}

Write $1=S0$ and use larger numerals as successor abbreviations.  Define
\[
 C(n,a,b)=I(n,a;u.b),\qquad
 Z(n)=C(n,1,0),\qquad
 a+b=I(b,a;u.Su).
\]
The update variable of $C$ is absent from its body.  Consequently
\[
 C(0,a,b)=a,\qquad C(Sn,a,b)=b.
\]
Numerical cases show that $Z(n)$ is $0$ or $1$.
The usual commutative-monoid equations for addition follow by literal
equation induction.  For example, $0+b=b$ and $Sa+b=S(a+b)$ are
proved by induction on $b$; these yield commutativity by a further
equation induction, and associativity follows by induction on the last
summand.  We use these equations to reorder finite sums.  The notation
$2a$ means $a+a$ and $3a$ means $a+a+a$; no multiplication operation
is needed for this bootstrap.

Put
\[
 N(n)=I(n,0;u.Z(u)),\qquad
 j(u)=Z(u)+2N(u),\qquad
 M(n)=I(n,0;u.j(u)).
\]
The range of $N$ is contained in $\{0,1\}$: its zero value is $0$,
and its successor value is a value of $Z$.  This uses numerical cases
on its counter, not induction on a disjunction.  Hence
\begin{equation}
 N(n)+Z(N(n))=1,\qquad
 Z(Z(N(n)))=N(n),\qquad N(n+2)=N(n).
 \label{eq:hybrid-parity}
\end{equation}
Here the last equation follows from two successor computations and
the preceding equation.  Literal-equation inductions give
\begin{equation}
 N(2n)=0,\qquad N(2n+1)=1,\qquad N(3n)=N(n).
 \label{eq:hybrid-parity-scaling}
\end{equation}
For the first equation the successor step adds two to the argument of
$N$; the second follows by one computation.  In the third, increasing
$n$ adds three to the left argument and therefore applies $Z$ once to
both sides of the induction equation.

Finite computations now give
\[
 j(0)=1,\qquad j(1)=2,\qquad j(2)=0.
\]
The proof below needs just these three values and
$M(Sn)=j(M(n))$; it does not assume a division algorithm or a theorem
characterizing $M$ as remainder modulo three.

\subsection{A ramp with a global scaling identity}

Define
\[
 G(n)=I(n,0;u.u+M(u)+1).
\]
The central invariant is
\begin{equation}
 M(G(n))=N(n).
 \label{eq:hybrid-ramp-invariant}
\end{equation}
Its zero case is immediate.  For its successor step, retain the
single induction equation and split the Boolean value $N(n)$.
If it is $0$, the next $G$ state is $G(n)+1$; its $M$ value is
$j(0)=1$, equal to $N(Sn)$.  If it is $1$, the next $G$ state
is $G(n)+2$; two computations of $M$ change its value from $1$
to $2$ and then to $0$, again equal to $N(Sn)$.  This proves
the invariant by literal-equation induction.

It follows that
\[
 G(Sn)=G(n)+N(n)+1,\qquad G(n+2)=G(n)+3.
\]
Equation induction therefore yields
\begin{equation}
 G(2n)=3n.
 \label{eq:hybrid-ramp-double}
\end{equation}
A separate literal-equation induction proves
\begin{equation}
 G(3n)=3G(n)+N(n).
 \label{eq:hybrid-ramp-triple}
\end{equation}
For completeness, its successor calculation is
\begin{align*}
 G(3Sn)
 &=G(3n)+N(3n)+4
   =3G(n)+2N(n)+4,\\
 3G(Sn)+N(Sn)
 &=3G(n)+3N(n)+3+Z(N(n)).
\end{align*}
The expressions agree in both Boolean cases for $N(n)$.
The base is $0=0$.  In particular,
\eqref{eq:hybrid-ramp-triple} was proved internally; it was not
inferred from an intended floor-function interpretation of $G$.

Set
\[
 H(u)=C(N(u),G(u),u).
\]
The crucial choice is to retain $u$ when its parity is $1$.
The preceding equations give
\begin{equation}
 H(2u)=3u,\qquad H(1)=1,\qquad H(3u)=3H(u).
 \label{eq:hybrid-factor-map}
\end{equation}
To check the last identity, use $N(3u)=N(u)$ and split the Boolean
$N(u)$.  At $0$, both sides are $3G(u)$ by
\eqref{eq:hybrid-ramp-triple}; at $1$, both are $3u$.
This argument does not require that an arbitrary number has already
been represented as $2v$ or $2v+1$.

\subsection{A comparator from scalar iteration}

For every fixed update $h$, equation induction on $n$, keeping $m,a$
and all parameters fixed, proves the semigroup identity
\[
 I(m+n,a;h)=I(n,I(m,a;h);h).
\]
The base is the zero computation and the step applies $h$ to the
induction equation. This identity will also be used in the later
pairing construction.


For a fixed update $h$, equation induction on $y$, with its initializer
$a$ and all parameters of $h$ fixed, proves the rotation identity
\begin{equation}
 I(Sy,a;h)=I(y,h(a);h).
 \label{eq:hybrid-rotation}
\end{equation}
Both sides have base $h(a)$, and their successor equations apply the
same $h$ to the previous equality.  Here and below $I(y,a;h)$ is
an abbreviation for one fixed finite scalar term, not quantification
over functions.

By \eqref{eq:hybrid-factor-map}, another equation induction gives
\begin{equation}
 I(y,3a;H)=3I(y,a;H).
 \label{eq:hybrid-iterator-scaling}
\end{equation}
Its base is $3a=3a$.  In the step the fixed-parameter induction
equation rewrites the new left side as
\[
 H(3I(y,a;H))=3H(I(y,a;H)).
\]
The initializer parameter stays $a$ throughout this induction.
Equation induction also gives $I(y,1;H)=1$, since $H(1)=1$.

Define
\[
 E(x)=I(x,1;u.2u),\qquad
 L(x,y)=N\bigl(I(y,E(x);H)\bigr).
\]
Then the following entire recursion specification is provable:
\begin{equation}
 L(0,y)=1,\qquad L(Sx,0)=0,\qquad L(Sx,Sy)=L(x,y).
 \label{eq:hybrid-comparator}
\end{equation}
The first equation uses $E(0)=1$ and the fixed-point identity for
$H$.  The second uses $E(Sx)=2E(x)$ and $N(2a)=0$.
For the third,
\begin{align*}
 L(Sx,Sy)
 &=N\bigl(I(Sy,2E(x);H)\bigr)\\
 &=N\bigl(I(y,H(2E(x));H)\bigr)\\
 &=N\bigl(I(y,3E(x);H)\bigr)\\
 &=N\bigl(3I(y,E(x);H)\bigr)\\
 &=N\bigl(I(y,E(x);H)\bigr)=L(x,y).
\end{align*}
The successive steps use rotation, $H(2a)=3a$, iterator scaling,
and $N(3a)=N(a)$.  Thus $L$ has the comparison equations at all
arguments internally.  Its standard-model meaning as the indicator
of $x\le y$ is motivation, not a premise of the following argument.

\subsection{Saturated iteration gives predecessor}

We now need only the three equations
\eqref{eq:hybrid-comparator}.  Define
\[
 B_m(u)=C(L(Su,m),u,Su),\qquad
 W(t,m)=I(t,0;u.B_m(u)).
\]
The subscript in $B_m$ records a fixed parameter; $B_m$ is an open
term abbreviation, not a function variable.

Numerical cases on $e$ prove the selector equation
\[
 C(e,Sa,Sb)=S C(e,a,b).
\]
Equivalently one may use equation induction on $e$, whose induction
hypothesis is unused.  Therefore
\begin{align}
 B_{Sm}(Su)
 &=C(L(SSu,Sm),Su,SSu)\notag\\
 &=C(L(Su,m),Su,SSu)=S B_m(u),
 \label{eq:hybrid-conjugacy}\\
 B_{Sm}(0)&=C(L(S0,Sm),0,S0)=C(L(0,m),0,S0)=S0.
 \label{eq:hybrid-conjugacy-base}
\end{align}

Literal-equation induction on $t$, with $m$ fixed, proves
\begin{equation}
 W(St,Sm)=S W(t,m).
 \label{eq:hybrid-saturated-shift}
\end{equation}
The base is \eqref{eq:hybrid-conjugacy-base}.  The successor step is
\begin{align*}
 W(SSt,Sm)
 &=B_{Sm}(W(St,Sm))\\
 &=B_{Sm}(S W(t,m))\\
 &=S B_m(W(t,m))=S W(St,m).
\end{align*}
Only the second line uses the current induction hypothesis, at the
same fixed $m$.

Two subsequent equation inductions give
\begin{equation}
 L(Sn,n)=0,\qquad W(n,n)=n.
 \label{eq:hybrid-diagonals}
\end{equation}
The first uses the comparator recurrence.  The second uses
\eqref{eq:hybrid-saturated-shift} in its successor step.  The
substitution $t=m=n$ is made into that already proved universally
parameterized equation; no parameter of an active induction hypothesis
is changed.  Consequently
\begin{equation}
 W(Sn,n)=B_n(n)=C(L(Sn,n),n,Sn)=n.
 \label{eq:hybrid-extra-step}
\end{equation}

Finally define the finite scalar term
\begin{equation}
 P(m)=I\bigl(m,0;u.C(L(SSu,m),u,Su)\bigr).
 \label{eq:hybrid-predecessor-term}
\end{equation}
Its zero equation is immediate.  To justify its successor equation
without assuming a separate extensionality rule for iterator bodies,
put
\[
 V(t,m)=I\bigl(t,0;u.C(L(SSu,Sm),u,Su)\bigr).
\]
The comparator recurrence identifies this update pointwise with
$B_m(u)$.  Equation induction on $t$, with $m$ fixed, gives
$V(t,m)=W(t,m)$: their base values are $0$, and the successor step
uses the pointwise update equation and then the induction equation.
Now \eqref{eq:hybrid-extra-step} yields
\[
 \boxed{\qquad P(0)=0,\qquad P(Sm)=V(Sm,m)=W(Sm,m)=m.\qquad}
\]

\subsection{Dependency boundary}

All terms in this section are built from $0,S,I$ by finite substitution.
All inductions have a single equation as their invariant.  Numerical
cases and equality transport are the already reconstructed logical
resources.  Successor injectivity may now be obtained by applying $P$
to an assumed equation $Sx=Sy$.  Equality tests, reflection of
quantifier-free formulas, pairing, and the recovery of general
primitive recursion can then be developed in the established order.
They were not used in order to manufacture predecessor.

This argument is a written internal derivation in the stated theory.
The finite implementation checks do not constitute a proof-assistant
formalization of this entire metatheorem.

% END


% BEGIN reference/equality_reflection.tex
% Standalone section.  Requires only amsmath and amssymb in the parent document.
% No theorem environments, macros, or bibliography keys are assumed.
% All arithmetic claims in this section are proved here.

\section{From literal-equation induction to full intuitionistic PRA}
\label{sec:eic-completion}\label{app:arithmetic}

This section supplies an unconditional completion theorem. It is adapted
with attribution from the current forward-left draft~\cite{forward}. In the
scalar application used here, every function occurring in the proof is
already a scalar term: predecessor is supplied by
Appendix~\ref{sec:hybrid-bootstrap}, the selector is
$C(n,a,b)=I(n,a;u.b)$, addition is $x+y=I(y,x;u.Su)$, and
truncated subtraction is $d(x,y)=I(y,x;u.Pu)$. All other functions below
are finite compositions of these. Thus this proof does not presuppose
that arbitrary primitive recursion has already been compiled. The QFI
conclusion first applies to the scalar language; the general PR compiler
comes afterward.  Literal-equation
induction, together with the primitive-recursive defining equations,
predecessor coverage, and separation of $0$ and $1$, proves induction for every
quantifier-free formula.  In particular, the proof does not assume a faithful
numerical equality test, stability or decidability of equality, successor
uniqueness, or any commutation law for an iterator.  We construct the equality
test and prove its reflection property for arbitrary elements of the
numerical domain.

\subsection{The exact hypotheses}

Work in ordinary intuitionistic first-order logic with equality and genuine
falsity $\bot$.  The language has one numerical sort and a symbol for each
fixed, finite primitive-recursive description.  Distinct descriptions are
not identified merely because they compute the same standard function.
Let $\mathsf{PRE}$ contain the defining equations for zero, successor,
projections, composition, and primitive recursion.  In particular, choose
descriptions $P$ and $C$ with equations
\begin{align}
 P(0)&=0, & P(Sx)&=x, \label{eq:eic-predecessor}\\
 C(0,a,b)&=a, & C(Sx,a,b)&=b. \label{eq:eic-selector}
\end{align}
Write $1=S0$.  In addition to $\mathsf{PRE}$, assume
\begin{align}
 \mathsf{Cov}:&\qquad
       \forall x\,\bigl(x=0\lor x=S(Px)\bigr),
       \label{eq:eic-coverage}\\
 \mathsf{Sep}:&\qquad 0=1\longrightarrow\bot,
       \label{eq:eic-separation}
\end{align}
and induction for each \emph{literal equation} $E(n,\vec p)$:
\begin{equation}
 \forall\vec p\,
 \left[
   \left(E(0,\vec p)\land
       \forall n\bigl(E(n,\vec p)\longrightarrow E(Sn,\vec p)\bigr)
   \right)
   \longrightarrow\forall n\,E(n,\vec p)
 \right].
 \label{eq:eic-induction}
\end{equation}
Here the schematic expression $E$ is a single equality of numerical
terms; the displayed universal quantifiers range only over numbers.
Every description and every instance of the schema is finite external
syntax.

The axiom formulation \eqref{eq:eic-induction} and the usual contextual
induction rule are interchangeable.  For example, if
\[
 \Gamma\vdash E(0,\vec p),
 \qquad
 \Gamma,E(z,\vec p)\vdash E(Sz,\vec p),
\]
where $z$ is absent from $\Gamma$ and the passive parameters, discharge the
second premise's equation and universally introduce $z$.  The axiom then
gives $\Gamma\vdash\forall n E(n,\vec p)$, which can be instantiated at any
target term.  Conversely, assume the antecedent of the axiom and apply the
contextual rule with fresh variables.  This uses the ordinary intuitionistic
logical rules and permits arbitrary, including quantified, formulas in
$\Gamma$.

\subsection{Constructor separation and two elementary tests}

Define the sign description and the countdown description by
\begin{equation}
 \operatorname{sg}(u)=C(u,0,1),
 \qquad
 K(u)=C(u,1,Pu).
 \label{eq:eic-sign-countdown}
\end{equation}
Their defining equations imply
\begin{align}
 \operatorname{sg}(0)&=0,
 &\operatorname{sg}(Sx)&=1,\label{eq:eic-sign-computation}\\
 K(0)&=1,
 &K(Sx)&=x.\label{eq:eic-countdown-computation}
\end{align}

\paragraph{Constructor laws.}
Successor is injective: applying $P$ to $Sx=Sy$ gives $x=y$ by
\eqref{eq:eic-predecessor}.  Successor never equals zero.  Indeed,
$Sx=0$ gives $1=0$ by applying $\operatorname{sg}$ and using
\eqref{eq:eic-sign-computation}; symmetry and
\eqref{eq:eic-separation} give $\bot$.  Consequently $1\ne0$, where
$A\ne B$ abbreviates $A=B\to\bot$.  None of these arguments uses
induction.

\paragraph{The sign test.}
For every $u$,
\begin{align}
 &\operatorname{sg}(u)=0\lor\operatorname{sg}(u)=1,
       \label{eq:eic-sign-range}\\
 &\operatorname{sg}(u)=0\quad\longleftrightarrow\quad u=0,
       \label{eq:eic-sign-zero}\\
 &\operatorname{sg}(u)=1\quad\longleftrightarrow\quad u\ne0.
       \label{eq:eic-sign-one}
\end{align}
To prove the range assertion, split \eqref{eq:eic-coverage}.  The branch
$u=0$ gives sign $0$, and the branch $u=S(Pu)$ gives sign $1$.
The implication $u=0\to\operatorname{sg}(u)=0$ is computation.
For its converse, the successor branch of coverage would make
$\operatorname{sg}(u)=0$ imply $1=0$, a contradiction; in that branch
explosion gives $u=0$ as well.  Similarly, if
$\operatorname{sg}(u)=1$ and $u=0$, computation gives $0=1$.
Conversely, $u\ne0$ excludes the zero branch of coverage, and the
successor branch computes the sign as $1$.

Coverage and \eqref{eq:eic-selector} also give, for all $u,v$,
\begin{equation}
 C(u,v,v)=v.
 \label{eq:eic-selector-idempotence}
\end{equation}
Both coverage branches compute to $v$, so this identity again needs no
induction.

\paragraph{The countdown test and the notation $W$.}
For every $u$,
\begin{equation}
 K(u)=0\quad\longleftrightarrow\quad u=1.
 \label{eq:eic-countdown-zero}
\end{equation}
If $u=1=S0$, \eqref{eq:eic-countdown-computation} gives $K(u)=0$.
Conversely, assume $K(u)=0$ and split coverage.  If $u=0$, then
$K(u)=1$, contradicting separation, and explosion gives $u=1$.
If $u=S(Pu)$, computation gives $K(u)=Pu$; hence $Pu=0$ and
$u=S0=1$.

If a presentation uses a description named $W$ with defining equations
\[
 W(0)=1,\qquad W(Sx)=x,
\]
then it proves $K(u)=W(u)$ for every $u$ already from coverage.
In the zero branch both sides are $1$; in the branch $u=S(Pu)$ both
sides are $Pu$.  Thus $K$ and $W$ are provably equal on the whole domain,
even if they are distinct descriptions.  This equality does not appeal to
standard-model agreement or to uniqueness of an arbitrary recurrence.

\subsection{Four elementary uses of equation induction}

Choose descriptions for addition and truncated subtraction, with recursion
on the second argument:
\begin{align}
 x+0&=x, &x+Sy&=S(x+y),\label{eq:eic-addition}\\
 d(x,0)&=x, &d(x,Sy)&=P(d(x,y)).\label{eq:eic-subtraction}
\end{align}
Using these equations and \eqref{eq:eic-induction}, we now derive
\begin{align}
 0+y&=y,\label{eq:eic-left-zero}\\
 (Sx)+y&=S(x+y),\label{eq:eic-left-successor}\\
 d(Sx,Sy)&=d(x,y),\label{eq:eic-cancel-successors}\\
 d(Sa,a)&=1.\label{eq:eic-diagonal}
\end{align}
Each induction below has a single equation as its induction formula.

For \eqref{eq:eic-left-zero}, induct on $y$.  The base is $0+0=0$.
The step is $0+Sy=S(0+y)=Sy$.

For \eqref{eq:eic-left-successor}, hold $x$ fixed and induct on $y$.
At zero both sides are $Sx$.  At a successor the induction hypothesis
gives
\[
 (Sx)+Sy=S\bigl((Sx)+y\bigr)
       =S\bigl(S(x+y)\bigr)=S(x+Sy).
\]
Combining \eqref{eq:eic-left-successor} with the defining equation for
addition gives the useful rearrangement
\begin{equation}
 (Sy)+r=y+Sr.
 \label{eq:eic-shift-successor}
\end{equation}
Indeed, both sides are $S(y+r)$.

For \eqref{eq:eic-cancel-successors}, hold $x$ fixed and induct on $y$.
The base computes as
\[
 d(Sx,S0)=P(d(Sx,0))=P(Sx)=x=d(x,0).
\]
Applying $P$ to the induction equation gives the step:
\[
 d(Sx,S(Sy))=P(d(Sx,Sy))=P(d(x,y))=d(x,Sy).
\]

Finally, induct on $a$ for \eqref{eq:eic-diagonal}.  The base is
$d(S0,0)=S0=1$.  At a successor,
\[
 d(S(Sa),Sa)=d(Sa,a)=1,
\]
using \eqref{eq:eic-cancel-successors} and the induction hypothesis.

\subsection{An equation invariant that reflects actual equality}

The remaining arithmetic ingredient is the following equation, valid for
all $x,y$:
\begin{equation}
 C\bigl(\operatorname{sg}(d(x,y)),\ x,\ y+d(x,y)\bigr)=x.
 \label{eq:eic-residual-invariant}
\end{equation}
Its meaning is that a nonzero subtraction residual reconstructs the
original argument by addition.  The selector expresses this assertion as
one literal equation, so equation induction applies directly.

Hold $x$ fixed and induct on $y$.  At $y=0$, subtraction gives residual
$x$, and \eqref{eq:eic-left-zero} gives $0+x=x$.  The desired equation
is consequently $C(\operatorname{sg}(x),x,x)=x$, which is
\eqref{eq:eic-selector-idempotence}.

For the step, assume \eqref{eq:eic-residual-invariant} at $y$ and
abbreviate
\[
 r=d(x,y),\qquad r'=Pr=d(x,Sy).
\]
These are local term abbreviations, not additional function symbols.
Split coverage for $r'$.

If $r'=0$, then $\operatorname{sg}(r')=0$, and the successor instance
of \eqref{eq:eic-residual-invariant} reduces immediately to $x=x$.

In the other branch, $r'=S(Pr')$, so constructor separation gives
$r'\ne0$.  Split coverage for $r$.  Its zero branch would imply
$r'=P0=0$, a contradiction, from which the desired step equation follows.
In its successor branch,
\[
 r=S(Pr)=Sr',\qquad \operatorname{sg}(r)=1.
\]
The induction equation therefore computes to $y+r=x$.  Meanwhile,
\eqref{eq:eic-shift-successor} gives
\[
 (Sy)+r'=y+Sr'=y+r=x.
\]
Since $r'$ is a successor, $\operatorname{sg}(r')=1$, so this is
exactly the branch selected in the successor instance of
\eqref{eq:eic-residual-invariant}.  This completes the equation induction.
All case distinctions came from coverage; no induction on an implication
or a disjunction was used.

\paragraph{Reflection at residual $1$.}
We have, for arbitrary $a,b$,
\begin{equation}
 d(Sa,b)=1\quad\longleftrightarrow\quad a=b.
 \label{eq:eic-residual-reflection}
\end{equation}
For the forward direction, substitute $x=Sa$, $y=b$ into
\eqref{eq:eic-residual-invariant}.  If $d(Sa,b)=1$, the sign is $1$
and the selector returns its second branch, giving
\[
 b+1=Sa.
\]
The defining equation for addition gives $b+1=S(b+0)=Sb$.  Thus
$Sb=Sa$, and applying $P$ gives $b=a$, hence $a=b$.
Conversely, $a=b$ and \eqref{eq:eic-diagonal} give
$d(Sa,b)=d(Sa,a)=1$.

Define the equality characteristic by the fixed description
\begin{equation}
 e(a,b)=\operatorname{sg}\bigl(K(d(Sa,b))\bigr).
 \label{eq:eic-equality-characteristic}
\end{equation}
The sign, countdown, and residual reflection lemmas give
\begin{align}
 e(a,b)=0&\quad\longleftrightarrow\quad a=b,
       \label{eq:eic-equality-reflection}\\
 e(a,b)=0&\ \lor\ e(a,b)=1.
       \label{eq:eic-equality-range}
\end{align}
In detail, $e(a,b)=0$ is equivalent to $K(d(Sa,b))=0$ by
\eqref{eq:eic-sign-zero}, which is equivalent to $d(Sa,b)=1$ by
\eqref{eq:eic-countdown-zero}, and then to $a=b$ by
\eqref{eq:eic-residual-reflection}.  The range assertion is
\eqref{eq:eic-sign-range} applied to $K(d(Sa,b))$.

The reflection property concerns the logic's actual equality throughout
the domain.  It was not assumed as part of a numerical comparison symbol,
and its proof did not restrict $a$ or $b$ to standard numerals.
It also yields decidable equality.  In the $e(a,b)=0$ branch,
\eqref{eq:eic-equality-reflection} gives $a=b$.  In the $e(a,b)=1$
branch, an assumption $a=b$ would give $e(a,b)=0$ and therefore $0=1$,
contradicting separation.  Hence
\begin{equation}
 \forall a\,\forall b\,\bigl(a=b\lor(a=b\to\bot)\bigr)
 \label{eq:eic-equality-decidable}
\end{equation}
is a theorem, not an extra hypothesis.

\subsection{Compilation of all quantifier-free formulas}

For each fixed quantifier-free formula $A(\vec p)$, construct a term
$b_A(\vec p)$ by external recursion on its finite syntax.  Take equality,
$\bot$, $\land$, $\lor$, and $\to$ as the primitive formula constructors;
negation is implication into $\bot$.  Set
\begin{align}
 b_{s=t}&=e(s,t), & b_{\bot}&=1,\label{eq:eic-boolean-atoms}\\
 b_{A\land D}&=C(b_A,b_D,1),\label{eq:eic-boolean-and}\\
 b_{A\lor D}&=C(b_A,0,b_D),\label{eq:eic-boolean-or}\\
 b_{A\to D}&=C(b_A,b_D,0).\label{eq:eic-boolean-implies}
\end{align}
Every term here is a finite composition of fixed primitive-recursive
descriptions; if desired, projections and composition give it a single
description symbol.  No internal quantification over formulas or programs
is involved.

An external induction on the syntax of $A$ proves, within the arithmetic
theory,
\begin{equation}
 \bigl(b_A=0\lor b_A=1\bigr)
 \quad\text{and}\quad
 \bigl(b_A=0\longleftrightarrow A\bigr).
 \label{eq:eic-boolean-correctness}
\end{equation}
Here are the full connective checks.  In every case $b_A=1$ implies
$\neg A$: if $A$ held, the already established implication
$A\to b_A=0$ would contradict separation.

\paragraph{Atomic formulas and falsity.}
For an equation, both claims are
\eqref{eq:eic-equality-reflection}--\eqref{eq:eic-equality-range}.
For falsity the range assertion is $1=0\lor1=1$, with its right disjunct
proved by reflexivity.  Also $1=0\to\bot$ follows from separation and
symmetry, while $\bot\to1=0$ follows by explosion.  Thus $1=0$ is
equivalent to $\bot$, as required.

\paragraph{Conjunction.}
Split $b_A=0\lor b_A=1$.  In the zero branch, $A$ holds and
$b_{A\land D}=b_D$.  Its range follows from the range of $b_D$,
and its being zero is equivalent to $D$, hence to $A\land D$ under
the known $A$.  In the one branch, $\neg A$ holds and
$b_{A\land D}=1$.  Its range is immediate.  If it were zero,
separation would give $\bot$, hence $A\land D$; conversely an
assumption $A\land D$ contradicts $\neg A$, and explosion gives
$b_{A\land D}=0$.

\paragraph{Disjunction.}
In the branch $b_A=0$, the formula $A$ holds and
$b_{A\lor D}=0$.  The range is immediate, and both $b_{A\lor D}=0$
and $A\lor D$ hold, proving their equivalence.  In the branch
$b_A=1$, we have $\neg A$ and $b_{A\lor D}=b_D$.  Its range follows
from the range of $b_D$.  If its value is zero, then $D$ and hence
$A\lor D$.  Conversely, from $A\lor D$, the $A$ alternative
contradicts $\neg A$ and yields $b_D=0$ by explosion, while the
$D$ alternative yields $b_D=0$ directly.  This proves the required
equivalence in both alternatives.

\paragraph{Implication.}
In the branch $b_A=0$, the formula $A$ holds and
$b_{A\to D}=b_D$.  Its range follows from the range of $b_D$.
If its value is zero, then $D$, which gives $A\to D$.  Conversely,
$A\to D$ and the known $A$ give $D$, hence $b_D=0$.
In the branch $b_A=1$, we have $\neg A$ and
$b_{A\to D}=0$.  Its range is immediate.  To prove $A\to D$,
assume $A$; this contradicts $\neg A$, and explosion gives $D$.
Thus both $b_{A\to D}=0$ and $A\to D$ hold in this branch, and
their equivalence follows.  This argument uses the proved two-value
range of the finite characteristic.  It does not use excluded middle
for arbitrary formulas or change the surrounding intuitionistic logic.

This completes the external syntactic induction proving
\eqref{eq:eic-boolean-correctness} for every quantifier-free formula.

\subsection{Recovery of induction and the full theory}

Let $A(n,\vec p)$ be any fixed quantifier-free formula.  Suppose an
arbitrary retained context $\Gamma$ proves its base and its induction
step:
\[
 \Gamma\vdash A(0,\vec p),
 \qquad
 \Gamma,A(z,\vec p)\vdash A(Sz,\vec p),
\]
with the usual freshness condition on $z$.  By
\eqref{eq:eic-boolean-correctness}, the base gives
$\Gamma\vdash b_A(0,\vec p)=0$.  Under the additional assumption
$b_A(z,\vec p)=0$, the same equivalence gives $A(z,\vec p)$,
the given step gives $A(Sz,\vec p)$, and the equivalence gives
$b_A(Sz,\vec p)=0$.  Apply literal-equation induction to
$b_A(n,\vec p)=0$ and then apply its reflection implication.
The result is $\Gamma\vdash A(t,\vec p)$ for every target term $t$.
Thus the theory proves the full contextual quantifier-free induction
rule, including in arbitrary quantified contexts.

\paragraph{Completion theorem.}
In the common language of fixed primitive-recursive descriptions, the
theory
\[
 \mathsf{PRE}+\mathsf{Cov}+\mathsf{Sep}+\mathsf{EI}
\]
has exactly the same sequents as ordinary first-order intuitionistic
primitive recursive arithmetic with all quantifier-free induction
instances.  The forward simulation follows from the recovered
quantifier-free induction, successor injectivity, and successor separation;
all primitive-recursive defining equations and all logical rules are
shared.  For the reverse simulation, literal-equation induction is an
instance of quantifier-free induction, and successor separation gives
$\mathsf{Sep}$.  Quantifier-free induction proves $\mathsf{Cov}$:
at zero choose its left disjunct, and at a successor use
$Sx=S(P(Sx))$, which follows from $P(Sx)=x$, for its right disjunct.

Both interpretations therefore use the full domain, defined by $x=x$,
and preserve every function symbol, equality, logical connective, and
quantifier.  Both composites have the identity isomorphism, with graph
$x=y$.  The theories are consequently identity-domain bi-interpretable
and equiconsistent.  This completion theorem is unconditional once the
literal-equation induction rule has been obtained; it has no dependency
on a classical successor-uniqueness reduction.

% END


% BEGIN reference/primitive_recursion.tex
\section{Pairing, primitive recursion, and internal numerical composites}
\label{sec:iterator-coding}
This appendix is adapted with attribution from the progress-closure
arithmetic development~\cite{progress}. Its prerequisites have now been
proved: scalar predecessor, arithmetic reflection, and full quantifier-free
induction. Pairing is a derived proof tool; it is not part of any checker.

\subsection{Pairing from a bounded square-root search}

Once QFI is available, a familiar square-based injection gives a
particularly direct compiler. Write $a\le b$ for $d(a,b)=0$.
The usual linear-order and semiring identities follow by fixed
quantifier-free inductions from the successor, addition, and monus
equations. Multiplication itself is the scalar term
\[
                        a\cdot b=I(b,0;u.u+a).
\]
Only primitive-recursive arithmetic is involved in the following
bounded root search; there is no minimization over an unbounded domain.

For a fixed parameter $n$, put
\[
 U_n(v)=C(d((Sv)^2,n),Sv,v),\qquad
 R_n(i)=I(i,0;v.U_n(v)),\qquad q(n)=R_n(n).
\]
The update increments $v$ exactly when $(v+1)^2\le n$, and
otherwise retains it. QFI in $i$, with $n$ fixed, proves
\[
 R_n(i)\le i,\qquad R_n(i)^2\le n,\qquad
 R_n(i)=i\ \vee\ n<(S(R_n(i)))^2.              \tag{Root invariant}
\]
The base has $R_n(0)=0$. If the update increments, its test
contradicts the strict-upper-bound alternative in the third invariant.
Thus the old state equals $i$, and the new state equals $Si$;
its square bound is the increment test. If the update does not increment, the
strict upper square bound holds and the state remains fixed.
An already established strict upper square bound forces every
subsequent update to remain fixed. These are exactly the cases
needed for the three displayed invariants.

At $i=n$, either the strict upper square bound already holds or
$R_n(n)=n$, in which case $n<(n+1)^2$ supplies it. Thus
\[
                        q(n)^2\le n<(q(n)+1)^2.           \tag{Root}
\]
Monotonicity of squares shows uniqueness of a number satisfying
Root. All formulas just used are quantifier free in the scalar
language, with fixed $n$ during the induction.

Now define
\[
 \pi(a,b)=(a+b)^2+a,\qquad
 L(n)=d(n,q(n)^2),\qquad H(n)=d(q(n),L(n)).
\]
For $k=a+b$, one has $k^2\le k^2+a<(k+1)^2$ since $a\le k$.
Root uniqueness gives $q(\pi(a,b))=a+b$, and therefore
\[
                    L(\pi(a,b))=a,\qquad H(\pi(a,b))=b.
\]
Surjectivity of $\pi$ is unnecessary. The original numerical domain
is never replaced by the range of this coding; pairing is used only
to store the finite state of a fixed primitive recursion. This
standard square pairing is a convenient reconstruction here, not
a claimed new arithmetic basis.

\subsection{Compiling each primitive recursion}

Suppose the fixed target description is
\[
 f(0,\vec p)=g(\vec p),\qquad
 f(Sn,\vec p)=h(n,f(n,\vec p),\vec p).
\]
Assume by external structural induction that the earlier $g,h$
already have scalar translations $g^*,h^*$. Define
\[
 W(u,\vec p)=\pi(SL(u),h^*(L(u),H(u),\vec p)),
\]
\[
 V(n,\vec p)=I(n,\pi(0,g^*(\vec p));u.W(u,\vec p)),
 \qquad f^*(n,\vec p)=H(V(n,\vec p)).
\]
The update has only its state variable bound, with every $p_i$
fixed. Equation induction proves $L(V(n,\vec p))=n$.
The two projection equations and scalar computation then prove
both defining equations for $f^*$ internally. Initial functions
and composition translate directly. This completes the translation
of every fixed finite PR description without adding a universal
evaluator.

\subsection{Why both numerical round trips are internal}

Two fixed scalar iterators with the same seed and provably equal
updates are internally equal at all counters: use literal-equation
induction comparing their values. Their base agrees, and their
successor equations plus congruence prove the step. The same
uniqueness argument works for ordinary primitive recursion.

For an ordinary PR description, induct externally on its finite
construction. The double translation satisfies its original
defining equations, so internal equation induction identifies it
with the original operation. For a native scalar description,
earlier translated updates agree by the external induction
hypothesis; scalar uniqueness identifies the double-translated
iterator. These proofs hold at arbitrary internal counters.

Thus the domain relation $x=y$ is provably a total bijection
preserving every fixed numerical operation in both composites.
The formula constructions and round trips in
Section~\ref{sec:choice-introduction} give the corresponding formula
composites. Both interpretations preserve falsity, taking native
$0=1$ to an ordinary contradiction and ordinary $\bot$ to native
$0=1$. Finite proof translations therefore establish equiconsistency.

% END


\clearpage
\begin{thebibliography}{99}
\raggedright
\bibitem{snapshot}
Vladimir Reshetnikov, repository maintainer.
\emph{ProveIt: current draft collection}, commit
\file{0300a290bedee5019c974291850b798378438b2e}.
\url{https://github.com/VladimirReshetnikov/ProveIt/tree/0300a290bedee5019c974291850b798378438b2e/docs/drafts}.
All archive references below refer to this immutable collection.

\bibitem{flat}
\emph{Flat computation for intuitionistic primitive recursive arithmetic}.
\file{ipra_flat_computation_2026-11-10.zip}, in~\cite{snapshot}.
Source of the null-tag uniform computation interface and the complete
attributed logical/arithmetic appendices used here.

\bibitem{constructed}
\emph{Constructing computation sources for iPRA}.
\file{ipra_constructed_computation_2026-11-10.zip}, in~\cite{snapshot}.

\bibitem{transient}
\emph{A one-state recursor for iPRA}.
\file{ipra_transient_seed_2026-11-10.zip}, in~\cite{snapshot}.
Source of the predecessor-packet joint profile, the 311-node comparator,
and a tied-seed bounding-coordinate obstruction.

\bibitem{counter}
\emph{Counter congruence suffices for iPRA}.
\file{ipra_counter_congruence_2026-11-10.zip}, in~\cite{snapshot}.
Source of the 672-byte complete control and related initializer obstructions.

\bibitem{scalar}
\emph{Scalar successor recursion for intuitionistic PRA}.
\file{ipra_scalar_successor_2026-11-10.zip}, in~\cite{snapshot}.

\bibitem{completion}
\emph{Completing the successor-seeded scalar basis for iPRA}.
\file{ipra_successor_seed_completion_2026-11-10.zip}, in~\cite{snapshot}.

\bibitem{choice}
\emph{One assumption at a time: a smaller iPRA proof kernel},
Choice-DI development, reproduced with attribution in the named-specialization,
generated-assumptions, and flat-computation archives of~\cite{snapshot}.

\bibitem{erasure}
Inherited scalar predecessor development, reproduced as
\file{reference/scalar_predecessor.tex} in~\cite{flat} and in this report.

\bibitem{forward}
Inherited forward-left equality-reflection development, reproduced as
\file{reference/equality_reflection.tex} in~\cite{flat} and in this report.

\bibitem{progress}
Inherited progress-closure pairing and primitive-recursion compiler,
reproduced as \file{reference/primitive_recursion.tex} in~\cite{flat}
and in this report.
\end{thebibliography}

% END

\end{document}
